\documentclass[open=any]{scrbook}
\usepackage{indentfirst}
\usepackage{amsmath,bm}
\usepackage{amssymb}
\usepackage{amsthm}
\usepackage{enumitem}
\usepackage{xcolor}
\usepackage{tikz-cd}
\usepackage{geometry}
\usepackage{tabto}
\geometry{
	a4paper,
	inner=20mm,
	outer=20mm,
	top=20mm,
	bottom=20mm,
}
\usepackage{hyperref}
\hypersetup{
	colorlinks=true,
	linkcolor=blue,
}
\usepackage[nameinlink]{cleveref}

\DeclareFontFamily{U}{mathc}{}
\DeclareFontShape{U}{mathc}{m}{it}{<-> mathc10}{}
\DeclareMathAlphabet{\mathcal}{U}{mathc}{m}{it}

% Create proposition/theorem style in main content
\newcommand{\CreateManualEnv}[2]{%
	% 1. Create the underlying amsthm environment
	\newtheorem{#1temp}{#2}
	% 2. Create custom wrapper environment dynamically
	\NewDocumentEnvironment{#1}{m o}{
		\expandafter\gdef\csname the#1temp\endcsname{##1}% % Sets \thetheoremtemp, etc.
		\IfValueTF{##2}{
			\begin{#1temp}[##2]
			}{
				\begin{#1temp}
				}
				\label{##1}
			}{
			\end{#1temp}
		}
	}
\CreateManualEnv{theorem}{Theorem}
\CreateManualEnv{proposition}{Proposition}
\CreateManualEnv{corollary}{Corollary}
\newtheorem*{proposition*}{Proposition}
\renewcommand{\proofname}{Sketch proof}

% Create definition style. Only for vertical separation.
\newtheoremstyle{noperiod}% name of the style
{\topsep}%             Space above
{\topsep}%             Space below
{\normalshape}%        Body font
{}%                    Indent amount
{}%                    Theorem head font (bold)
{}%                    Punctuation after theorem head (LEAVE EMPTY to remove the period)
{\parindent}%          Space after theorem head
{}%                    Theorem head spec (can be left empty)
\theoremstyle{noperiod}
\newtheorem*{definition}{}

% Create exercise style
\theoremstyle{definition}
\newtheorem{exercise}{}[section]

\crefname{exercise}{exercise}{exercises}

\begin{document}
	{
		\hypersetup{
			colorlinks=true,
			linkcolor=black,
		}
		\tableofcontents
	}
	\clearpage
	\chapter{Sheaves}

\section*{PRESHEAVES}

\begin{definition}
	Let $X$ be a topological space. A \textcolor{orange}{presheaf} of sets is a functor $\mathcal{F}:\mathbf{Open}(X)^\text{op}\rightarrow\mathbf{Set}$, where $\operatorname{ob}\left(\mathbf{Open}(X)\right)$ is the set of all open subsets of $X$, and for any $U, V \in \operatorname{ob}\left(\mathbf{Open}(X)\right)$,
\begin{equation*}
	\operatorname{Hom}\left(U,V\right)=\left\{\begin{aligned}
		&\left\{\iota_{VU}\right\}, && \text{if }V\subseteq U;\\
		&\varnothing, && \text{if }V\nsubseteq U.
	\end{aligned}\right.
\end{equation*}
where $\iota_{UV}$ denotes the unique inclusion morphism from $V$ to $U$.

The category $\mathbf{Set}$ in the definition of a presheaf can be replaced by an arbitrary category $\mathcal{C}$, and the corresponding functor is called a \textcolor{orange}{$\mathcal{C}$-valued presheaf}. 
\end{definition}

\begin{definition}
	An element $s \in \mathcal{F}(U)$ is called a \textcolor{orange}{section} of $\mathcal{F}$ over $U$.
	
	The set $\mathcal{F}(U)$ is called the \textcolor{orange}{set of sections} (or simply the \textcolor{orange}{sections}) over $U$.
	
	The map $\mathcal{F}\!\left(\iota_{VU}\right)$ is called the \textcolor{orange}{restriction map}, and $\mathcal{F}\!\left(\iota_{VU}\right)(s)$ is written to be $s|_V$ if $s\in\mathcal{F}(U)$.
\end{definition}

\begin{definition}
	Let $\mathcal{F}$ and $\mathcal{G}$ be presheaves on $X$. A \textcolor{orange}{morphism} $\varphi:\mathcal{F}\rightarrow\mathcal{G}$ is a collection of maps $\varphi_U:\mathcal{F}(U)\rightarrow\mathcal{G}(U)$ for each open set $U \subseteq X$, such that whenever $V\subseteq U$, the diagram
	\begin{equation*}
		\begin{tikzcd}
			\mathcal{F}(U)\arrow[r,"\mathcal{F}\left(\iota_{VU}\right)"]\arrow[d,"\varphi_U"]&[2em]\mathcal{F}(V)\arrow[d,"\varphi_V"]\\
			\mathcal{G}(U)\arrow[r,"\mathcal{G}\left(\iota_{VU}\right)"]&\mathcal{G}(V)
		\end{tikzcd}
	\end{equation*}
	commutes. All presheaves on $X$ form a category, denoted by \textcolor{orange}{$\mathbf{PSh}(X)$}.
	
	An isomorphism is defined to be a morphism having a two-sided inverse.
\end{definition}

\begin{definition}
	Let $\varphi:\mathcal{F}\rightarrow\mathcal{G}$ be a morphism of presheaves. The \textcolor{orange}{presheaf kernel} of $\varphi$, \textcolor{orange}{presheaf cokernel} of $\varphi$, and \textcolor{orange}{presheaf image} of $\varphi$ are defined to be the presheaves given by $U\mapsto\ker(\varphi_U)$, $U\mapsto\operatorname{coker}(\varphi_U)$ and $U\mapsto\operatorname{im}(\varphi_U)$, respectively.
\end{definition}

\begin{definition}
	Let $\mathcal{F}$ be a presheaf on $X$. The \textcolor{orange}{stalk} of $\mathcal{F}$ at $x \in X$ is the direct limit
	\begin{equation*}
		\mathcal{F}_x=\varinjlim_{U\ni x}\mathcal{F}(U).
	\end{equation*}
	
	The elements of a stalk $\mathcal{F}_x$ are the \textcolor{orange}{germs} of sections of $\mathcal{F}$ at point $x$.
	
	Denote $\mathcolor{orange}{s_x}\in\mathcal{F}_x$ to be the image of $s\in\mathcal{F}(U)$ under the canonical morphism of the direct limit.
\end{definition}

\begin{definition}
	The \textcolor{orange}{\'etal\'e space} of the presheaf $\mathcal{F}$ is the disjoint union of all stalks
	\begin{equation*}
		E(\mathcal{F})=\coprod_{x\in X}\mathcal{F}_x
	\end{equation*}
	equipped with the topology generated by the basis of sets $U_s=\left\{s_x\in\mathcal{F}_x\mid x\in U\right\}$ for all open sets $U\subseteq X$ and all sections $s\in\mathcal{F}(U)$, where $s_x$ denotes the image of $s$ in the stalk $\mathcal{F}_x$.
	
	There is a canonical projection $\pi: E(\mathcal{F}) \rightarrow X$ defined by $s_x \mapsto x$ for each $s_x \in \mathcal{F}_x$. The topology of the \'etal\'e space ensures that $\pi$ is a continuous map (more precisely, a local homeomorphism).
\end{definition}


\section*{SHEAVES}

\begin{definition}
	A presheaf $\mathcal{F}$ on a topological space $X$ is a \textcolor{orange}{sheaf} if, for any open set $U\subseteq X$ and any open covering $\left\{V_i\right\}_{i\in I}$ of $U$, the following conditions hold:
	\begin{enumerate}[label=\textup{(\arabic*)}]
		\item \textcolor{red}{Locality:} For any sections $s, t\in\mathcal{F}(U)$, if $s|_{V_i}=t|_{V_i}$ for all $i\in I$, then $s = t$.
		\item \textcolor{red}{Gluing:} For a family of sections $\left\{s_i\in\mathcal{F}\!\left(V_i\right)\right\}_{i\in I}$ satisfying $s_i|_{V_i\cap V_j}=s_j|_{V_i\cap V_j}$ for all $i, j\in I$, there exists a section $s\in\mathcal{F}(U)$ such that $s|_{V_i}=s_i$ for all $i\in I$.
	\end{enumerate}
	
	The category of sheaves on $X$, denoted by \textcolor{orange}{$\mathbf{Sh}(X)$}, is defined as the full subcategory of $\mathbf{PSh}(X)$ consisting of those presheaves that satisfy both the locality and gluing properties.
\end{definition}

\begin{proposition}{1.1}
	Let $\varphi:\mathcal{F}\rightarrow\mathcal{G}$ be a morphism of presheaves on a topological space $X$. If $\varphi$ is an isomorphism, then the induced map $\varphi_x:\mathcal{F}_x\rightarrow\mathcal{G}_x$ is an isomorphism for every $x\in X$. If in addition $\mathcal{F}$ and $\mathcal{G}$ are sheaves, then $\varphi$ is an isomorphism iff $\varphi_x$ is an isomorphism for all $x\in X$.
\end{proposition}

\begin{definition}
	Let $\mathcal{F}$ be a presheaf. The \textcolor{orange}{associated sheaf} (or \textcolor{orange}{sheafification}) of $\mathcal{F}$ is a sheaf $\mathcal{F}^+$ which maps each open set $U$ to the set of maps $s:U\rightarrow E(\mathcal{F})$, such that
	\begin{enumerate}[label=\textup{(\arabic*)}]
		\item for each $x\in U$, $s(x)\in\mathcal{F}_x$;
		\item for each $x\in U$, there exists an open neighborhood $V\subseteq U$ of $x$ and a section $t\in\mathcal{F}(V)$ such that $t_y=s(y)$ for all $y\in V$.
	\end{enumerate}
	
	There is a canonical morphism of presheaves $\theta:\mathcal{F}\rightarrow\mathcal{F}^+$ defined by $\mathcal{F}(U)\mapsto\mathcal{F}^+(U)$.
\end{definition}

\begin{proposition}{1.2}[Universal property of sheafification]
	Let $\theta:\mathcal{F}\rightarrow\mathcal{F}^+$ be the canonical morphism. Then for any sheaf $\mathcal{G}$ and any morphism $\varphi:\mathcal{F} \rightarrow\mathcal{G}$, there exists a unique morphism of presheaves $\psi:\mathcal{F}^+\rightarrow\mathcal{G}$ such that $\varphi=\psi\circ\theta$.
\end{proposition}

\begin{definition}
	Let $\varphi:\mathcal{F}\rightarrow\mathcal{G}$ be a morphism of sheaves. The \textcolor{orange}{kernel} of $\varphi$ is defined to be $\ker_\mathbf{PSh}\varphi$; the \textcolor{orange}{cokernel} and \textcolor{orange}{image} of $\varphi$ are defined to be $(\operatorname{coker}_\mathbf{PSh}\varphi)^+$ and $(\operatorname{im}_\mathbf{PSh}\varphi)^+$, respectively.
	
	The morphism $\varphi$ is \textcolor{orange}{injective} if $\ker\varphi=0$. It is \textcolor{orange}{surjective} if $\operatorname{im}\varphi=\mathcal{G}$.
	
	Let $\mathcal{F}^\prime$ be a subsheaf of $\mathcal{F}$. The \textcolor{orange}{quotient sheaf} $\mathcal{F}/\mathcal{F}^\prime$ is defined to be the associated sheaf of the presheaf $U\mapsto\mathcal{F}(U)/\mathcal{F}^\prime(U)$. Use the definition of associated sheaf, it follows that $(\mathcal{F}/\mathcal{F}^\prime)_x=\mathcal{F}_x/\mathcal{F}_x^\prime$.
\end{definition}

\begin{definition}
	Let $f:X\rightarrow Y$ be a continuous map of topological spaces.
	
	The \textcolor{orange}{direct image} of a sheaf $\mathcal{F}$ on $X$ is the sheaf $f_*\mathcal{F}$ on $Y$ defined by $f_*\mathcal{F}(V)=\mathcal{F}\!\left(f^{-1}(V)\right)$.
	
	The \textcolor{orange}{inverse image} of a sheaf $\mathcal{G}$ on $Y$ is the sheaf $f^{-1}\mathcal{G}$ on $X$ defined by
	\begin{equation*}
		f^{-1}\mathcal{G}(U)=\left(\varinjlim_{V\supseteq f(U)}\mathcal{G}(V)\right)^+.
	\end{equation*}
	
	The \'etal\'e space of $f^{-1}\mathcal{G}$ is the pullback $E(\mathcal{G})\times_YX$ satisfying the commutative diagram
	\begin{equation*}
		\begin{tikzcd}
			E(\mathcal{G})\times_YX\arrow[r]\arrow[d]&E(\mathcal{G})\arrow[d,"\pi_Y"]\\
			X\arrow[r,"f"]&Y
		\end{tikzcd}
	\end{equation*}
\end{definition}

\begin{definition}
	Let $\mathcal{F}$ be a sheaf on $X$ and $Z$ is a subset of $X$. The \textcolor{orange}{restriction} of $\mathcal{F}$ to $Z$ is the inverse image $i^{-1}\mathcal{F}$, where $i:Z\rightarrow X$ is the inclusion map, denote by \textcolor{orange}{$\mathcal{F}|_Z$}.
\end{definition}
	\newpage
	\chapter{Schemes}

\section*{RINGED SPACES}

\begin{definition}
	A \textcolor{orange}{ringed space} is a pair $\left(X,\mathcal{O}_X\right)$ consisting of a topological space $X$ and a sheaf of rings $\mathcal{O}_X:\mathbf{Open}(X)^\text{op}\rightarrow\mathbf{Ring}$. The sheaf $\mathcal{O}_X$ is called the \textcolor{orange}{structure sheaf} of $\left(X,\mathcal{O}_X\right)$.
	
	A ringed space $\left(X,\mathcal{O}_X\right)$ is called a \textcolor{orange}{locally ringed space} if the stalk $\mathcal{O}_{X,x}$ is a local ring for every point $x\in X$.
\end{definition}

\begin{definition}
	A \textcolor{orange}{morphism of ringed spaces} from $\left(X,\mathcal{O}_X\right)$ to $\left(Y,\mathcal{O}_Y\right)$ is a pair $\left(f,f^\#\right)$ consisting of a continuous map $f:X\rightarrow Y$ and a map $f^\#:\mathcal{O}_Y\rightarrow f_*\mathcal{O}_X$ of sheaves of rings on $Y$.
	
	A homomorphism of local rings $\varphi:A\rightarrow B$ is called a \textcolor{orange}{local homomorphism} if $\varphi^{-1}\!\left(\mathfrak{m}_B\right)=\mathfrak{m}_A$.
	
	A \textcolor{orange}{morphism of locally ringed spaces} is a morphism $\left(f,f^\#\right)$ of ringed spaces such that each induced homomorphism $f^\#_x:\mathcal{O}_{Y,f(x)}\rightarrow\mathcal{O}_{X,x}$ of local rings is a local homomorphism.
	
	An isomorphism is defined to be a morphism having a two-sided inverse.
\end{definition}

For any ring $R$, $\operatorname{Spec}(R)$ is a ringed space (\textcolor{blue}{Atiyah-Macdonald Exercise 3.24}). Consider the inclusion map $\mathbb{Z}_{(p)}\hookrightarrow\mathbb{Q}$ which induces a morphism of ringed spaces $f:\operatorname{Spec}(\mathbb{Q})\rightarrow\operatorname{Spec}\!\left(\mathbb{Z}_{(p)}\right)$. However, there is another morphism $g:\operatorname{Spec}(\mathbb{Q})\rightarrow\operatorname{Spec}\!\left(\mathbb{Z}_{(p)}\right)$ defined by $(0)\mapsto(p)$. To exclude such examples, we must restrict our attention to morphisms of locally ringed spaces.

\section*{SCHEMES}

\begin{definition}
	An \textcolor{orange}{affine scheme} is a locally ringed space $\left(X,\mathcal{O}_X\right)$ that is isomorphic (as a locally ringed space) to $\left(\operatorname{Spec}(R),\mathcal{O}_{\operatorname{Spec}(R)}\right)$ for some ring $R$.
	
	A \textcolor{orange}{scheme} is a locally ringed space $\left(X,\mathcal{O}_X\right)$ admitting an open cover $\left\{U_i\right\}$ of $X$ such that each open subscheme $\left(U_i,\mathcal{O}_X|_{U_i}\right)$ is an affine scheme.
	
	A \textcolor{orange}{morphism of schemes} is a morphism as locally ringed spaces.

	Denote \textcolor{orange}{$\operatorname{sp}(X)$} to be the underlying topological space of a scheme $\left(X,\mathcal{O}_X\right)$.
\end{definition}

Let $S=\bigoplus_{n\geqslant0}S_n$ be a graded ring. Define \textcolor{orange}{$\operatorname{Proj}(S)$} to be the set of all homogeneous prime ideals of $S$ that do not contain the irrelevant ideal $S_+=\bigoplus_{n>0}S_n$. Define a topology on $\operatorname{Proj}(S)$ by taking the closed subsets to be of the form $V(\mathfrak{a})$ for any homogeneous ideal $\mathfrak{a}$.

\begin{proposition}{2.5}
	For any homogeneous element $f\in S_+$, let $D_+(f)=\left\{\mathfrak{p}\in\operatorname{Proj}(S)\mid f\notin\mathfrak{p}\right\}$ be the basic open set, and let $\mathcal{O}_{\operatorname{Proj}(S)}\!\left(D_+(f)\right)=S_{(f)}$ where $S_{(f)}$ is the degree 0 part of the localization $S_f$ ($x/f^k\in S_{(f)}$ means $\deg x=k\deg f$). Then $\left(\operatorname{Proj}(S),\mathcal{O}_{\operatorname{Proj}(S)}\right)$ is a scheme.
\end{proposition}

\begin{definition}
	Let $S$ be a fixed scheme. A \textcolor{orange}{scheme over $S$} is a scheme $X$, together with a morphism $X\rightarrow S$.
	
	If $X$ and $Y$ are schemes over $S$, an \textcolor{orange}{$S$-morphism} of $X$ to $Y$ is a morphism $f:X\rightarrow Y$ such that the diagram
	\begin{equation*}
		\begin{tikzcd}
			X\arrow[r,"f"]\arrow[d]&Y\arrow[ld]\\
			S
		\end{tikzcd}
	\end{equation*}
	commutes.
	
	Denote \textcolor{orange}{$\mathbf{Sch}$} to be the category of schemes, and \textcolor{orange}{$\mathbf{Sch}(S)$} to be the category of schemes over $S$.
	
	If $R$ is a ring, denote \textcolor{orange}{$\mathbf{Sch}(R)$} to be the category of schemes over $\operatorname{Spec}(R)$.
\end{definition}

\begin{proposition}{2.6}
	Let $k$ be an algebraically closed field. There is a natural fully faithful functor $t:\mathbf{Var}(k)\rightarrow\mathbf{Sch}(k)$ from the category of varieties over $k$ to schemes over $k$. For any variety $V$, its topological space is homeomorphic to the set of closed points of $\operatorname{sp}(t(V))$, and its sheaf of regular functions is obtained by restricting the structure sheaf of $t(V)$ via this homeomorphism.
\end{proposition}

\section*{DEFINITIONS AND PROPOSITIONS FROM EXERCISES}

\begin{definition}
	Let $X$ be scheme. The \textcolor{orange}{residue field} of $x$ on $X$ is the field $k(x)=\mathcal{O}_{X,x}/\mathfrak{m}_x$.
\end{definition}
	\newpage
	\chapter{First Properties of Schemes}

\section*{SCHEMES}

\begin{definition}
	A scheme is \textcolor{orange}{connected} if its topological space is connected.
	
	In particular, an affine scheme $\operatorname{Spec}R$ is connected iff $R$ has no idepotent elements.
\end{definition}

\begin{definition}
	A scheme $X$ is \textcolor{orange}{reduced} if for every open set $U$, the ring $\mathcal{O}_X(U)$ has no nilpotent elements.
	
	A scheme is \textcolor{orange}{irreducible} if its topological space is irreducible.
	
	A scheme $X$ is \textcolor{orange}{integral} if for every open set $U\subseteq X$, the ring $\mathcal{O}_X(U)$ is an integral domain.
\end{definition}

\begin{proposition}{3.1}
	A scheme is integral iff it is both reduced and irreducible.
\end{proposition}

\begin{definition}
	A scheme is \textcolor{orange}{locally Noetherian} if it can be covered by open affine subset $\operatorname{Spec}\!\left(A_i\right)$, where each $A_i$ is Noetherian.
	
	A scheme is \textcolor{orange}{Noetherian} if it is locally Noetherian and quasi-compact.
\end{definition}

\begin{proposition}{3.2}
	A scheme is locally Noetherian iff for every open affine subset $U=\operatorname{Spec}A$, $A$ is a Noetherian ring.
\end{proposition}

\section*{MORPHISMS}

\begin{definition}
	A morphism $f:X\rightarrow Y$ is \textcolor{orange}{locally of finite type}, if there exists a covering of $Y$ by open affine subsets $V_i=\operatorname{Spec}\!\left(B_i\right)$, such that for each $i$, $f^{-1}\!\left(V_i\right)$ can be covered by open affine subsets $U_{ij}=\operatorname{Spec}\!\left(A_{ij}\right)$, where each $A_{ij}$ is a finitely generated $B_i$-algebra.
	
	The morphism $f$ is of \textcolor{orange}{finite type} if in addition for each $i$, the cover $\left\{U_{ij}\right\}$ has a finite subcover.

\end{definition}

\begin{definition}
	A morphism $f:X\rightarrow Y$ is \textcolor{orange}{finite}, if there exists a covering of $Y$ by open affine subsets $V_i=\operatorname{Spec\!}\left(B_i\right)$, such that for each $i$, $f^{-1}\!\left(V_i\right)$ can be convered by open affine subsets $U_{ij}=\operatorname{Spec}\!\left(A_{ij}\right)$, where each $A_{ij}$ is a finitely generated $B_i$-module.
	
	{
	\setlist[enumerate]{itemsep=0pt,topsep=0pt,parsep=0pt}
	A morphism $f:X\rightarrow Y$ is \textcolor{orange}{quasi-finite}, if \tabto{9cm}\textcolor{magenta}{Inequivalent definitions}
	\begin{enumerate}[label=\textup{(\roman*)},ref=(\roman*)]
		\item\label{quasifinite1} for every point $y\in Y$, $f^{-1}(y)$ is a finite set. \tabto{8cm}\textcolor{magenta}{Hartshorne}
		\item\label{quasifinite2} \labelcref{quasifinite1} plus $f^{-1}(y)$ is a discrete set. \tabto{8cm}\textcolor{magenta}{Grothendieck SGA}
		\item\label{quasifinite3} \labelcref{quasifinite2} plus $f$ is of finite type. \tabto{8cm}\textcolor{magenta}{Grothendieck EGA}
	\end{enumerate}
	}
\end{definition}

\begin{definition}
	An \textcolor{orange}{open subscheme} of a scheme $X$ is an open subset $U$ with structure sheaf $\mathcal{O}_X|_U$.
	
	An \textcolor{orange}{open immersion} is a morphism $f:X\rightarrow Y$ which induces an isomorphism of $X$ with an open subscheme of $Y$.
\end{definition}

\begin{definition}
	A \textcolor{orange}{closed immersion} is a morphism $f:Y\rightarrow X$ such that $f$ induces a homeomorphism from $\operatorname{sp}(Y)$ to a closed subset of $\operatorname{sp}(X)$, and the induced map $f^\#:\mathcal{O}_X\rightarrow f_*\mathcal{O}_Y$ is surjective.
	
	A \textcolor{orange}{closed subscheme} of a scheme $X$ is an equivalence class of closed immersions, where $f:Y\rightarrow X$ and $f^\prime:Y^\prime\rightarrow X$ are equivalent if there is an isomorphism $i:Y^\prime\rightarrow Y$ such that $f^\prime=f\circ i$.
\end{definition}

\section*{DIMENSION AND PRODUCTS}

\begin{definition}
	The \textcolor{orange}{dimension} of a scheme $X$ is its dimension as a topological space.
	
	If $Z$ is an irreducible closed subset of $X$, the \textcolor{orange}{codimension} $\operatorname{codim}(Z,X)$ is the supremum of integers $n$ such that there exists a chain
	\begin{equation*}
		Z=Z_0<Z_1<\cdots<Z_n
	\end{equation*}
	of distinct closed irreducible subsets of $X$. If $Y$ is any closed subset of $X$, define
	\begin{equation*}
		\operatorname{codim}(Y,X)=\inf_{Z\subseteq Y}\operatorname{codim}(Z,X),
	\end{equation*}
	where the infimum is taken over all closed irreducible subsets of $Y$.
\end{definition}

\begin{definition}
	Let $S$ be a scheme, and let $X$, $Y$ be schemes over $S$. The \textcolor{orange}{fiber product} of $X$ and $Y$ over $S$ is the categorical product of $X$ and $Y$ in $\mathbf{Sch}(S)$, equivalently, the pullback of the morphism $X\rightarrow S$ and $Y\rightarrow S$ in $\mathbf{Sch}$. We denote the fiber product by \textcolor{orange}{$X\times_SY$}.
\end{definition}


\section*{DEFINITIONS AND PROPOSITIONS FROM EXERCISES}

\begin{definition}
	A morphism $f:X\rightarrow Y$ is \textcolor{orange}{quasi-compact} if the inverse image of open quasi-compact subset of $Y$ is quasi-compact on $X$.
\end{definition}
	\newpage
	\chapter{Separated and Proper Morphisms}

\section*{SEPARATED MORPHISMS}

\begin{definition}
	Let $f:X\rightarrow Y$ be the morphism of schemes. The \textcolor{orange}{diagonal morphism} is the unique morphism $\Delta:X\rightarrow X\times_YX$ satisfying $\pi_k\circ\Delta=\mathrm{id}_X$ for $k=1,2$, where $\pi_k:X\times_YX\rightarrow X$ is the canonical projection to the $k$-th component.
\end{definition}

\begin{definition}
	A morphism $f:X\rightarrow Y$ is \textcolor{orange}{separated} if the diagonal morphism $\Delta$ is a closed immersion.
	
	A morphism $f:X\rightarrow Y$ is \textcolor{orange}{quasi-separated} if the diagonal morphism $\Delta$ is quasi-compact. 
\end{definition}
	
\begin{definition}
	A scheme $X$ is \textcolor{orange}{separated over $Y$} if there exists a separated morphism $f:X\rightarrow Y$. A scheme $X$ is said to be \textcolor{orange}{separated} if it is separated over $\operatorname{Spec}\mathbb{Z}$.
\end{definition}

\begin{proposition}{4.1}
	Any morphism between affine schemes is separated.
\end{proposition}

\begin{corollary}{4.2}
	A morphism $f:X\rightarrow Y$ is separated iff the image of the diagonal morphism is a closed subset of $X\times_YX$.
\end{corollary}

\begin{theorem}{4.3}[\textcolor{red}{Valuative criterion of separatedness}]
	Let $f:X\rightarrow Y$ be a morphism of schemes. Then the following are equivalent:
	\begin{enumerate}[label=\textup{(\roman*)}]
		\item $f$ is separated;
		\item $f$ is quasi-separated, and for every valuation ring $R$ with fraction field $K$, every morphism $\operatorname{Spec}R\rightarrow Y$, and every morphism $\operatorname{Spec}K\rightarrow X$, there exists at most one morphism $h:\operatorname{Spec}R\rightarrow X$ such that the diagram
		\begin{equation*}
			\begin{tikzcd}
				\operatorname{Spec}K\arrow[r]\arrow[d,"i"]&X\arrow[d,"f"]\\
				\operatorname{Spec}R\arrow[ur,"h"]\arrow[r]&Y
			\end{tikzcd}
		\end{equation*}
		commutes.
	\end{enumerate}
\end{theorem}

\section*{PROPER MORPHISMS}

\begin{definition}
	A morphism $f:X\rightarrow Y$ \textcolor{orange}{universally closed} if for every scheme $Z$ with a morphism $Z\rightarrow Y$, the projection from the fiber product $X\times_YZ\rightarrow Z$ is a closed map of the underlying topological spaces.
\end{definition}

\begin{definition}
	A morphism $f:X\rightarrow Y$ is \textcolor{orange}{proper} if it is separated, of finite type, and universally closed.
\end{definition}

\begin{theorem}{4.7}[\textcolor{red}{Valuative criterion of properness}]
	Let $f:X\rightarrow Y$ be a morphism of schemes. Then the following are equivalent:
	\begin{enumerate}[label=\textup{(\roman*)}]
		\item $f$ is proper;
		\item $f$ is quasi separated and of finite type, and for every valuation ring $R$ with fraction field $K$, every morphism $\operatorname{Spec}R\rightarrow Y$, and every morphism $\operatorname{Spec}K\rightarrow X$, there exists a unique morphism $h:\operatorname{Spec}R\rightarrow X$ such that the diagram
		\begin{equation*}
			\begin{tikzcd}
				\operatorname{Spec}K\arrow[r]\arrow[d,"i"]&X\arrow[d,"f"]\\
				\operatorname{Spec}R\arrow[ur,"h"]\arrow[r]&Y
			\end{tikzcd}
		\end{equation*}
		commutes.
	\end{enumerate}
\end{theorem}

\begin{definition}
	Let $Y$ be a scheme. The \textcolor{orange}{projective $n$-space over $Y$} is $\mathbb{P}_\mathbb{Z}^n\times_{\operatorname{Spec}\mathbb{Z}}Y$, denoted by \textcolor{orange}{$\mathbb{P}_Y^n$}.
	
	A morphism $f:X\rightarrow Y$ of schemes is \textcolor{orange}{projective} if it factors into $X\xrightarrow{i}\mathbb{P}_Y^n\xrightarrow{\pi}Y$, where $i$ is a closed immersion and $\pi$ is the canonical projection.
	
	A morphism $f:X\rightarrow Y$ of schemes is \textcolor{orange}{quasi-projective} if it factors into $X\xrightarrow{j}X^\prime\xrightarrow{g}Y$, where $j$ is an open immersion and $g$ is a projective morphism.
\end{definition}

\begin{theorem}{4.9}
	A projective morphism of Noetherian schemes is proper. A quasi-projective morphism of Noetherian schemes is of finite type and separated.
\end{theorem}

\begin{proposition}{4.10}
	Let $k$ be an algebraically closed field. The image of the functor $t:\mathbf{Var}(k)\rightarrow\mathbf{Sch}(k)$ is the set of quasi-projective integral schemes over $k$. The image of the set of projective varieties is the set of projective integral schemes.
\end{proposition}

\begin{definition}
	An \textcolor{orange}{abstract variety} is an integral separated scheme of finite type over an algebraically closed field $k$.
	
	A \textcolor{orange}{complete variety} is an abstract variety which is proper over the field $k$.
\end{definition}

\section*{DEFINITIONS AND PROPOSITIONS FROM EXERCISES}

	\newpage
	\chapter{Sheaves of Modules}

\section*{SHEAVES OF MODULES}

\begin{definition}
	Let $\left(X,\mathcal{O}_X\right)$ be a ringed space. A \textcolor{orange}{sheaf of $\mathcal{O}_X$-modules} is a sheaf $\mathcal{F}:\mathbf{Open}(X)^\text{op}\rightarrow\mathbf{Ab}$, such that for each open set $U\subseteq X$, the group $\mathcal{F}(U)$ is an $\mathcal{O}_X(U)$-module, and for each inclusion of open sets $V\subseteq U$, the restriction homomorphism $\mathcal{F}(U)\rightarrow\mathcal{F}(V)$ is compatible with the module structures via the ring homomorphism $\mathcal{O}_X(U)\rightarrow\mathcal{O}_X(V)$.
	
	A \textcolor{orange}{morphism of sheaves of $\mathcal{O}_X$-modules} is a morphism $\mathcal{F}\rightarrow\mathcal{G}$ of sheaves, such that for each open set $U\subseteq X$, the map $\mathcal{F}(U)\rightarrow\mathcal{G}(U)$ is a homomorphism of $\mathcal{O}_X(U)$-modules.
\end{definition}

\begin{definition}
	Let $\mathcal{F}$ and $\mathcal{G}$ are $\mathcal{O}_X$-modules. The presheaf
	\begin{equation*}
		U\mapsto\operatorname{Hom}_{\mathcal{O}_X|_U}\!\left(\mathcal{F}|_U,\mathcal{G}|_U\right)
	\end{equation*}
	is a sheaf, which is called the \textcolor{orange}{sheaf $\mathcal{Hom}$}, and denote by \textcolor{orange}{$\mathcal{Hom}_{\mathcal{O}_X}(\mathcal{F},\mathcal{G})$}.
	
	The \textcolor{orange}{tensor product} of $\mathcal{F}$ and $\mathcal{G}$ is the sheaf associated to the preshheaf
	\begin{equation*}
		U\mapsto\mathcal{F}(U)\otimes_{\mathcal{O}_X(U)}\mathcal{G}(U),
	\end{equation*}
	denote by \textcolor{orange}{$\mathcal{F}\otimes_{\mathcal{O}_X}\mathcal{G}$}.
\end{definition}

\begin{definition}
	An $\mathcal{O}_X$-module $\mathcal{F}$ is \textcolor{orange}{free} if it is isomorphic to a direct sum of copies of $\mathcal{O}_X$.
	
	An $\mathcal{O}_X$-module $\mathcal{F}$ is \textcolor{orange}{locally free} if $X$ can be covered by open sets $U$ for which $\mathcal{F}|_U$ is a free $\mathcal{O}_X|_U$-module.
	
	The \textcolor{orange}{rank} of $\mathcal{F}$ on an open set $U$ is the number $r_U$ such that $\mathcal{F}|_U\cong\mathcal{O}_U^{\oplus r_U}$.
	
	A locally free sheaf of rank 1 is called an \textcolor{orange}{invertible sheaf}.
\end{definition}

\begin{definition}
	A \textcolor{orange}{sheaf of ideals} on $X$ is a sheaf of modules $\mathcal{I}$ which is a subsheaf of $\mathcal{O}_X$. In other words, for every open set $U$, $\mathcal{I}(U)$ is an ideal in $\mathcal{O}_X(U)$.
\end{definition}

\begin{definition}
	Let $f:\left(X,\mathcal{O}_X\right)\rightarrow\left(Y,\mathcal{O}_Y\right)$ be a morphism of ringed spaces, $\mathcal{F}$ be an $\mathcal{O}_X$-module and $\mathcal{G}$ be an $\mathcal{O}_Y$-module.	The \textcolor{orange}{direct image} of $\mathcal{F}$ is the sheaf $f_*\mathcal{F}$. The \textcolor{orange}{inverse image} of $\mathcal{G}$ is the sheaf $f^{-1}\mathcal{G}\otimes_{f^{-1}\mathcal{O}_Y}\mathcal{O}_X$, denote by \textcolor{orange}{$f^*\mathcal{G}$}.
\end{definition}

\section*{COHERENT SHEAVES}

\begin{definition}
	Let $A$ be a ring and $M$ be an $A$-module. The \textcolor{orange}{associated sheaf} of $M$ on $\operatorname{Spec}A$ is a sheaf $\tilde{M}$ which maps each open set $U$ to the set of maps $s:U\rightarrow\coprod_{\mathfrak{p}\in U}M_\mathfrak{p}$, such that
	\begin{enumerate}[label=\textup{(\arabic*)}]
		\item for each $\mathfrak{p}\in U$, $s(\mathfrak{p})\in M_\mathfrak{p}$;
		\item for each $\mathfrak{p}\in U$, there exists an open neighborhood $V\subseteq U$ of $\mathfrak{p}$, elements $m\in M$ and $f\in A$, such that $m/f=s(\mathfrak{q})$ for all $\mathfrak{q}\in V$ satisfying $f\notin\mathfrak{q}$.
	\end{enumerate}
\end{definition}

\begin{definition}
	Let $\left(X,\mathcal{O}_X\right)$ be a scheme. A sheaf of $\mathcal{O}_X$-modules $\mathcal{F}$ is \textcolor{orange}{quasi-coherent} if $X$ can be covered by open affine subsets $U_i=\operatorname{Spec}\!\left(A_i\right)$, such that for each $i$ there is an $A_i$-module $M_i$ with $\mathcal{F}|_{U_i}\cong\tilde{M_i}$. We said that $\mathcal{F}$ is \textcolor{orange}{coherent} if furthermore each $M_i$ can be taken to be a finitely generated $A_i$-module.
\end{definition}

\begin{definition}
	Let $A$ be a ring and $M$ be an $A$-module. The \textcolor{orange}{associated sheaf} of $M$ on $\operatorname{Spec}A$ is a sheaf $\tilde{M}$ which maps each open set $U$ to the set of maps $s:U\rightarrow\coprod_{\mathfrak{p}\in U}M_\mathfrak{p}$, such that
	\begin{enumerate}[label=\textup{(\arabic*)}]
		\item for each $\mathfrak{p}\in U$, $s(\mathfrak{p})\in M_\mathfrak{p}$;
		\item for each $\mathfrak{p}\in U$, there exists an open neighborhood $V\subseteq U$ of $\mathfrak{p}$, elements $m\in M$ and $f\in A$, such that $m/f=s(\mathfrak{q})$ for all $\mathfrak{q}\in V$ satisfying $f\notin\mathfrak{q}$.
	\end{enumerate}
\end{definition}


\section*{TWISTED SHEAVES}

\begin{definition}
	Let $S$ be a graded ring and $M$ be a graded $S$-module. The \textcolor{orange}{associated sheaf} of $M$ on $\operatorname{Proj}S$ is a sheaf $\tilde{M}$ which maps each open set $U$ to the set of maps $s:U\rightarrow\coprod_{\mathfrak{p}\in U}M_{(\mathfrak{p})}$, such that
	\begin{enumerate}[label=\textup{(\arabic*)}]
		\item for each $\mathfrak{p}\in U$, $s(\mathfrak{p})\in M_\mathfrak{p}$;
		\item for each $\mathfrak{p}\in U$, there exists an open neighborhood $V\subseteq U$ of $\mathfrak{p}$, elements $m\in M$ and $f\in A$, such that $m/f=s(\mathfrak{q})$ for all $\mathfrak{q}\in V$ satisfying $f\notin\mathfrak{q}$.
	\end{enumerate}
\end{definition}

\begin{definition}
	Let $S$ be a graded ring, and let $X=\operatorname{Proj}S$. For any $n\in\mathbb{Z}$, define $\mathcolor{orange}{\widetilde{S(n)}}=\mathcal{O}_X(n)$. $\mathcal{O}_X(1)$ is called the \textcolor{orange}{twisting sheaf} of Serre. For any sheaf of $\mathcal{O}_X$ moduels $\mathcal{F}$, the \textcolor{orange}{twisted sheaf} $\mathcal{F}\otimes_{\mathcal{O}_X}\mathcal{O}_X(n)$ is denoted by \textcolor{orange}{$\mathcal{F}(n)$}.
\end{definition}


	\newpage
	\chapter{Chain Conditions}

Let $\Sigma$ be a partially ordered set. $\Sigma$ said to be satisfying the \textcolor{orange}{ascending chain condition} (\textcolor{orange}{a.c.c}) if, for every increasing sequence $x_1\leqslant x_2\leqslant\cdots$ in $\Sigma$ there exists $n$ such that $x_n=x_{n+1}=\cdots$.

Similarly, $\Sigma$ said to be satisfying the \textcolor{orange}{descending chain condition} (\textcolor{orange}{d.c.c}) if, for every decreasing sequence $x_1\geqslant x_2\geqslant\cdots$ in $\Sigma$ there exists $n$ such that $x_n=x_{n+1}=\cdots$.

A module satisfying the ascending chain condition is said to be \textcolor{orange}{Noetherian}. A module satisfying the descending chain condition is said to be \textcolor{orange}{Artinian}.

\begin{proposition}{6.2}
	$M$ is a Noetherian $A$-module $\Leftrightarrow$ every submodule of $M$ is finitely generated.
\end{proposition}

\begin{proposition}{6.3}
	Let $0\rightarrow M^\prime\rightarrow M\rightarrow M^{\prime\prime}\rightarrow0$ be an exact sequence of $A$-modules. Then
	\begin{enumerate}[label=\textup{(\roman*)}]
		\item $M$ is Noetherian $\Leftrightarrow$ $M^\prime$ and $M^{\prime\prime}$ are Noetherian.
		\item $M$ is Artinian $\Leftrightarrow$ $M^\prime$ and $M^{\prime\prime}$ are Artinian.
	\end{enumerate}
\end{proposition}

\begin{proposition}{6.5}
	Let $A$ be a Noetherian (resp. Artinian) ring, $M$ a finitely generated $A$-module. Then $M$ is Noetherian (resp. Artinian).
\end{proposition}

\begin{definition}
	A \textcolor{orange}{chain} of submodules of a module $M$ is a sequence $\left(M_i\right)$ of submodules of $M$ such that
	\begin{equation*}
		M=M_0\supsetneq M_1\supsetneq\cdots\supsetneq M_n=0.
	\end{equation*}
	The \textcolor{orange}{length} of the chain is $n$. A \textcolor{orange}{composition series} of $M$ is a maximal chain, that is one which no extra submodules can be inserted, or equivalently $M_{i-1}/M_i$ is simple.
\end{definition}

\begin{proposition}{6.7}
	Suppose that $M$ has a composition series of length $n$. Then every composition series of $M$ has length $n$, and every chain can be extended to a composition series.
\end{proposition}

The length of a composition series of $M$ is called the \textcolor{orange}{length} of $M$, denoted as $l(M)$.

\begin{proposition}{6.10}
	For a vector space, 
	\begin{center}
		finite dimension $\Longleftrightarrow$ finite length $\Longleftrightarrow$ a.c.c. $\Longleftrightarrow$ d.c.c..
	\end{center}
	If the conditions are satisfied, length equals dimension.
\end{proposition}

\section*{DEFINITIONS AND PROPOSITIONS FROM EXERCISES}

A topological space $X$ is said to be \textcolor{orange}{Noetherian} if the open subsets of $X$ satisfy the ascending chain condition. Or equivalently, the closed subsets satisfy descending chain condition.
	\newpage
	\chapter{Noetherian Rings}

\begin{proposition}{7.1}
	Any homomorphic image of Noetherian ring is Noetherian.
\end{proposition}

\begin{proposition}{7.2}
	Let $A$ be a subring of $B$; soppose that $A$ is Noetherian and that $B$ is finitely generated as an $A$-module. Then $B$ is Noetherian.
\end{proposition}

\begin{proposition}{7.3}
	If $A$ is Noetherian and $S$ is any multiplicatively closed subset of $A$, then $S^{-1}A$ is Noetherian.
\end{proposition}

\begin{theorem}{7.5}[\textcolor{red}{Hilbert's basis theorem}]
	If $A$ is Noetherian, then the polynomial ring $A[x]$ is Noetherian.
\end{theorem}

\begin{proposition}{7.8}
	Let $A\subseteq B\subseteq C$ be rings. Suppose that $A$ is Noetherian, that $C$ is finitely generated as an $A$-algebra and that $C$ is either (i) finitely generated as a $B$-module or (ii) integral over $B$. Then $B$ is finitely generated as an $A$-algebra.
\end{proposition}

\begin{proposition}{7.9}[\textcolor{red}{Zariski's lemma}]
	Let $k$ be a field, $E$ a finitely generated $k$-algebra. If $E$ is a field then it is a finite algebraic extension of $k$.
\end{proposition}

\section*{PRIMARY DECOMPOSITION IN NOETHERIAN RINGS}

\begin{proposition}{7.11}
	In a Noetherian ring every ideal is a finite intersection of irreducible ideals.
\end{proposition}

\begin{proposition}{7.12}
	In a Noetherian ring every irreducible ideal is primary.
\end{proposition}

\begin{proposition}{7.13}
	In a Noetherian ring every ideal has a primary decomposition.
\end{proposition}

\begin{proposition}{7.14}
	In a Noetherian ring every ideal is contains a power of its radical.
\end{proposition}

\begin{corollary}{7.15}
	In a Noetherian ring the nilradical is nilpotent.
\end{corollary}

\begin{corollary}{7.16}
	Let $A$ be a Noetherian ring, $m$ a maximal ideal of $A$, $\mathfrak{q}$ any ideal of $A$. Then the following are equivalent:
	\begin{enumerate}[label=\textup{(\roman*)},ref=\thecorollarytemp(\roman*)]\crefalias{enumi}{corollary}
		\item\label{7.16.1} $\mathfrak{q}$ is $\mathfrak{m}$-primary;
		\item\label{7.16.2} $\sqrt{\mathfrak{q}}=\mathfrak{m}$;
		\item\label{7.16.3} $\mathfrak{m}^n\subseteq\mathfrak{q}\subseteq\mathfrak{m}$ for some $n>0$.
	\end{enumerate}
\end{corollary}

\begin{proposition}{7.17}
	Let $\mathfrak{a}\neq(1)$ be an ideal in a Noetherian ring. Then the prime ideals which belong to $\mathfrak{a}$ are precisely the prime ideals which occor in the set of ideals $(\mathfrak{a}:x)$.
\end{proposition}

\section*{DEFINITIONS AND PROPOSITIONS FROM EXERCISES}

\begin{definition}
	A subset of a topological space $X$ is \textcolor{orange}{constructible set} if it is a finite union of locally closed sets.
\end{definition}
	\newpage
	\chapter{Artinian Rings}

\begin{proposition}{8.1}
	In an Artinian ring every prime ideal is maximal.
\end{proposition}

\begin{proposition}{8.3}
	An Artinian ring has only a finite number of maximal ideals.
\end{proposition}

\begin{proposition}{8.4}
	In an Artinian ring the nilradical is nilpotent.
\end{proposition}

\begin{definition}
	The \textcolor{orange}{dimension} of $A$ is the supremum of the length of all chains of prime ideals in $A$.
\end{definition}

\begin{proposition}{8.5}
	A ring $A$ is Artinian $\Leftrightarrow$ $A$ is Noetherian and $\operatorname{dim}(A)=0$.
\end{proposition}

\begin{proposition}{8.6}
	Let $A$ be a Noetherian local ring, $\mathfrak{m}$ its maximal ideal. Then exactly one of the following statements is true:
	\begin{enumerate}[label=\textup{(\roman*)}]
		\item $\mathfrak{m}^n\neq\mathfrak{m}^{n+1}$ for all $n$;
		\item $\mathfrak{m}^n=0$ for some $n$. In this case $A$ is an Artinian local ring.
	\end{enumerate}
\end{proposition}

\begin{proposition}{8.7}[Structure theorem of Artinian rings]
	An Artinaian ring $A$ is uniquely (up to isomorphism) a finite direct product of Artinian local rings. 
\end{proposition}

\begin{proposition}{8.8}
	Let $A$ be an Artinian local ring. Then the following are equivalent:
	\begin{enumerate}[label=\textup{(\roman*)}]
		\item every ideal in $A$ is principal;
		\item the maximal ideal $\mathfrak{m}$ is principal;
		\item $\operatorname{dim}_{A/\mathfrak{m}}\left(\mathfrak{m}\middle/\mathfrak{m}^2\right)\leqslant1$.
	\end{enumerate}
\end{proposition}
	\newpage
	\chapter{Discrete Valuation Rings and Dedekind Domain}

\begin{proposition}{9.1}
	Let $A$ be a Noetherian domain of dimension 1. Then every non-zero ideal $\mathfrak{a}$ in $A$ can be uniquely expressed as a product of primary ideals whose radicals are all distinct.
\end{proposition}

\section*{DISCRETE VALUATION RINGS}

Let $K$ be a field. A \textcolor{orange}{discrete valuation} on $K$ is a mapping $v$ of $K^*$ onto $\mathbb{Z}$ such that
\begin{enumerate}[label=\textup{(\arabic*)}]
	\item $v{xy}=v(x)+v(y)$;
	\item $v(x+y)\geqslant\min\left(v(x),v(y)\right)$.
\end{enumerate}
The set consisting of all $x\in K^*$ such that $v(x)\geqslant0$ called the \textcolor{orange}{valuation ring} of $v$.

An integral domain $A$ is a \textcolor{orange}{discrete valuation ring} if there is a discrete valuation $v$ of its field of fractions $K$ such that $A$ is the valuation ring of $v$.

\begin{proposition}{9.2}
	Let $A$ be a Noetherian local domain of dimension 1, $\mathfrak{m}$ its maximal ideal. Then the following are equivalent:
	\begin{enumerate}[label=\textup{(\roman*)},ref=\thepropositiontemp(\roman*)]\crefalias{enumi}{proposition}
		\item\label{9.2.1} $A$ is a discrete valuation ring;
		\item\label{9.2.2} $A$ is integrally closed;
		\item\label{9.2.3} $\mathfrak{m}$ is a principal ideal;
		\item\label{9.2.4} $\dim_k\left(\mathfrak{m}/\mathfrak{m}^2\right)=1$;
		\item\label{9.2.5} Every non-zero ideal is a power of $\mathfrak{m}$;
		\item\label{9.2.6} There exists $x\in A$ such that every non-zero ideal is of the form $\left(x^k\right)$ for $k\geqslant0$.
	\end{enumerate}
\end{proposition}

\section*{DEDEKIND DOMAINS}

\begin{theorem}{9.3}
	Let $A$ be a Noetherian local domain of dimension 1. Then the following are equivalent:
	\begin{enumerate}[label=\textup{(\roman*)},ref=\thetheoremtemp(\roman*)]\crefalias{enumi}{theorem}
		\item\label{9.3.1} $A$ is integrally closed;
		\item\label{9.3.2} Every primary ideal in $A$ is a prime power;
		\item\label{9.3.3} Every local ring $A_\mathfrak{p}$ with $\mathfrak{p}\neq0$ is a discrete valuation ring.
	\end{enumerate}
	A ring satisfying the conditions is called a \textcolor{orange}{Dedekind domain}.
\end{theorem}

\section*{FRACTIONAL IDEALS}

Let $A$ be an integral domain, $K$ its field of fractions. An $A$-submodule $M$ of $K$ is a \textcolor{orange}{fractional ideal} of $A$ if $xM\subseteq A$ for some $x\neq0$.

An $A$-submodule $M$ of $K$ is an \textcolor{orange}{invertible ideal} if there exists a submodule $N$ of $K$ such that $MN=A$. Then $N=(A:M)$, and $M$ is finitely generated fractional ideal.

\begin{proposition}{9.6}
	For a fractional ideal $M$, the following are equivalent:
	\begin{enumerate}[label=\textup{(\roman*)}]
		\item $M$ is invertible;
		\item $M$ is finitely generated and , for each prime ideal $\mathfrak{p}$, $M_\mathfrak{p}$ is invertible;
		\item $M$ is finitely generated and , for each maximal ideal $\mathfrak{m}$, $M_\mathfrak{m}$ is invertible;
	\end{enumerate}
\end{proposition}

\begin{theorem}{9.8}
	Let $A$ be an integral domain. Then $A$ is a Dedekind domain $\Leftrightarrow$ every non-zero fractional ideal of $A$ is invertible.
\end{theorem}

	\newpage
	% Exercises section
	\chapter*{EXERCISES}
	\addcontentsline{toc}{chapter}{EXERCISES}
	\setlist[enumerate]{align=left,leftmargin=0pt,itemindent=!,itemsep=0pt,topsep=0pt,parsep=0pt}
	\setlist[itemize]{itemsep=0pt,topsep=0pt,parsep=0pt}
	\markboth{EXERCISES}{EXERCISES}
	\renewcommand{\thesection}{\arabic{section}}
	%\everymath{\color{red!50!black}}
	%\AtBeginEnvironment{equation*}{\color{red!50!black}}
	%\tikzset{every picture/.style={red!50!black}}
	\section{Sheaves}

\begin{exercise}\label{ex1.1}
	Apply the definition on $\mathcal{F}(U)$ and $\mathcal{F}\!\left(\iota_{VU}\right)$.
\end{exercise}

\begin{exercise}\label{ex1.2}
	\begin{enumerate}[label=(\alph*),ref=\theexercise(\alph*)]\crefalias{enumi}{exercise}
		\item\label{ex1.2.1} $\operatorname{im}\!\left(\varphi_x\right)=\operatorname{im}\!\left(\mathcal{F}_x\rightarrow\mathcal{G}_x\right)$, and $(\operatorname{im}\varphi)_x=\displaystyle\varinjlim_{U\ni x}\operatorname{im}\!\left(\mathcal{F}(U)\rightarrow\mathcal{G}(U)\right)$. The diagram
		\begin{equation*}
			\begin{tikzcd}
				\mathcal{F}(U)\arrow[r]\arrow[d,"\varphi_U"]&\mathcal{F}_x\arrow[d,"\varphi_x"]\\
				\mathcal{G}(U)\arrow[r]&\mathcal{G}_x
			\end{tikzcd}
			\qquad\text{can be refined to}\qquad
			\begin{tikzcd}
				\mathcal{F}(U)\arrow[r]\arrow[d,"\varphi_U"]&\mathcal{F}_x\arrow[d]\\
				\operatorname{im}\!\left(\varphi_U\right)\arrow[r]\arrow[d,hook]&\displaystyle\varinjlim_{U\ni x}\operatorname{im}\!\left(\varphi_U\right)\arrow[d]\\
				\mathcal{G}(U)\arrow[r]&\mathcal{G}_x
			\end{tikzcd}.
		\end{equation*}
		Use \textcolor{cyan}{Atiyah-Macdonald Exercise 2.15(i)} to prove $\displaystyle\varinjlim_{U\ni x}\operatorname{im}\!\left(\varphi_U\right)\subseteq\operatorname{im}\!\left(\varphi_x\right)$ and $\displaystyle\varinjlim_{U\ni x}\operatorname{im}\!\left(\varphi_U\right)\supseteq\operatorname{im}\!\left(\varphi_x\right)$. Use the \textcolor{cyan}{exactness of direct limits} to prove the $\ker$ case. \textcolor{red}{This is also correct for presheaves.}
		\item\label{ex1.2.2} $\hookrightarrow$: Use \hyperref[ex1.2.1]{(a)} to reduce to $(\ker\varphi)_x=0$ $\Rightarrow$ $\ker\varphi=0$. Let $s\in\mathcal{F}(U)$. For each $x\in U$, $s_x=0$, where $s_x\in\mathcal{F}_x$. By \textcolor{cyan}{Atiyah-Macdonald Exercise 2.15(ii)}, there exists $V_x\ni x$ such that $s|_{V_x}=0$. When $x$ runs over all points of $U$ we have an open cover $\left\{V_x\right\}$. Then use the locality axiom.
		
		$\twoheadrightarrow$: Use \hyperref[ex1.2.1]{(a)} to reduce to $(\operatorname{im}\varphi)_x=\mathcal{G}_x$ $\Rightarrow$ $\operatorname{im}\varphi=\mathcal{G}$. Let $t\in\mathcal{G}(U)$. For each $x\in U$, $t_x\in\mathcal{G}_x=(\operatorname{im}\varphi)_x$. By \textcolor{cyan}{Atiyah-Macdonald Exercise 2.15(i)}, there exist $V_x\ni x$ and $s_x\in(\operatorname{im}\varphi)\!\left(V_x\right)$ such that $\left(s_x\right)_x=t_x$. Since $\left(t|_{V_x}\right)_x=t_x$, \textcolor{cyan}{Atiyah-Macdonald Exercise 2.15(i)} again gives $W_x\subseteq V_x$, $x\in W_x$ such that $s_x|_{W_x}=t|_{W_x}$. Thus, we have a family of sections $\left\{s_x|_{W_x}\in(\operatorname{im}\varphi)\!\left(W_x\right)\right\}$ which satisfies $s_x|_{W_x\cap W_y}=t|_{W_x\cap W_y}=s_y|_{W_x\cap W_y}$. Then use the gluing axiom.
		
		\textcolor{magenta}{Note that the equation $s_x|_{W_x\cap W_y}=s_y|_{W_x\cap W_y}$ uses the injectiveness of $\operatorname{im}\varphi\hookrightarrow\mathcal{G}$ implicitly.}
		
		\textcolor{green!70!black}{Another approach is to use \Cref{1.1} on $(\ker\varphi)_x=0$ and $(\operatorname{im}\varphi)_x=\mathcal{G}_x$.}
	\end{enumerate}
\end{exercise}

\begin{exercise}\label{ex1.3}
	\begin{enumerate}[label=(\alph*),ref=\theexercise(\alph*)]\crefalias{enumi}{exercise}
		\item\label{ex1.3.1} $\Rightarrow$: Use gluing axiom. $\Leftarrow$: Use \Cref{ex1.2.1}.
		\item\label{ex1.3.2} Let $\varphi:\mathcal{O}_\mathbb{C}\rightarrow\mathcal{O}_{\mathbb{C}^\times}$, $f\mapsto\mathrm{e}^f$ be the sheaf morphism of holomorphic functions. Take $U=\mathbb{C}^\times$, then $\varphi_U:\mathcal{O}_\mathbb{C}(U)\rightarrow\mathcal{O}_{\mathbb{C}^\times}(U)$ cannot be surjective.
	\end{enumerate}
\end{exercise}

\begin{exercise}\label{ex1.4}
	\begin{enumerate}[label=(\alph*),ref=\theexercise(\alph*)]\crefalias{enumi}{exercise}
		\item\label{ex1.4.1} If $\varphi^+\!\left(s_1\right)=\varphi^+\!\left(s_2\right)$, $\varphi_x^+\!\left(s_{1,x}\right)=\varphi_x^+\!\left(s_{2,x}\right)$. We also have $\mathcal{F}_x\cong\mathcal{F}_x^+$ by the definition of sheafification. Then use \Cref{ex1.2.2}.
		\item \textcolor{magenta}{$\operatorname{im}\varphi$ is a presheaf but may not be a sheaf.} By \hyperref[ex1.4.1]{(a)}, $(\operatorname{im}\varphi)^+\rightarrow\mathcal{G}^+=\mathcal{G}$ is injective, so $(\operatorname{im}\varphi)^+$ can be identified as a subsheaf of $\mathcal{G}$.
	\end{enumerate}
\end{exercise}

\begin{exercise}\label{ex1.5}
	\Cref{1.1} and \Cref{ex1.2.2}.
\end{exercise}

\begin{exercise}\label{ex1.6}
	Work on stalks. Use \Cref{ex1.2.1} to find the kernel. \textcolor{green!70!black}{Or use \Cref{1.2}.}
\end{exercise}

\begin{exercise}\label{ex1.7}
	Use \Cref{ex1.2.1} to identify the maps on stalks, then use \Cref{1.1}. \textcolor{green!70!black}{Or use \Cref{1.2}.}
\end{exercise}

\begin{exercise}\label{ex1.8}
	Denote $0\rightarrow\mathcal{F}^\prime\xrightarrow{\varphi}\mathcal{F}\xrightarrow{\psi}\mathcal{F}^{\prime\prime}$. Use \hyperref[ex1.2.2]{\Cref*{ex1.2.2}$\hookrightarrow$} for injectivity of $\varphi_U$.
	
	Let $s\in\mathcal{F}(U)$ such that $\psi_U(s)=0$. By \hyperref[ex1.2]{\Cref{ex1.2}(c)}, there is $s^\prime\in\mathcal{F}^\prime(U)$ such that $\varphi_x(s^\prime_x)=s_x$. By \textcolor{cyan}{Atiyah-Macdonald Exercise 2.15} and similar to \hyperref[ex1.2.2]{\Cref*{ex1.2.2}$\twoheadrightarrow$}, there is $t_x\in\mathcal{F}\!\left(V_x\right)$ such that $\varphi_{V_x}\!\left(t_x\right)=s|_{V_x}$. Therefore, $\varphi_{V_x\cap V_y}\!\left(t_x|_{V_x\cap V_y}\right)=s|_{V_x\cap V_y}=\varphi_{V_x\cap V_y}\!\left(t_y|_{V_x\cap V_y}\right)$. Due to $\varphi_{V_x\cap V_y}$ is injective, using the gluing axiom of $\mathcal{F}^\prime$, we have $\varphi_{V_x}\left(t|_{V_x}\right)=s|_{V_x}$ for some $t\in U$. Since $\varphi$ is compatible with restriction maps, $\varphi_{V_x}\!\left(t|_{V_x}\right)=\left(\varphi_U(t)\right)\!|_{V_x}$. By the locality axiom of $\mathcal{F}$, $\varphi_U(t)=s$.
\end{exercise}

\begin{exercise}\label{ex1.9}
	Define the canonical projection and the canonical injection.
\end{exercise}

\begin{exercise}\label{ex1.10}
	$\varphi_U:\varinjlim\mathcal{F}_i(U)\rightarrow\mathcal{G}(U)$ is unique for each $U$. Follow \textcolor{red}{red} $\rightarrow$ \textcolor{orange}{orange} $\rightarrow$ \textcolor{green}{green} $\rightarrow$ \textcolor{cyan}{cyan} $\rightarrow$ \textcolor{blue}{blue} in the left diagram to prove $\varphi$ is compatible with the restriction maps, so $\varphi$ is a morphism of presheaves.
	\begin{equation*}
		\begin{tikzcd}
			\mathcal{F}_i(U)\arrow[rrr,shift right=.2ex,orange]\arrow[rrr,shift left=.2ex,green]\arrow[dr,shift right=.2ex,blue]\arrow[dr,shift left=.2ex,red]\arrow[ddr,cyan]&&&\mathcal{F}_i(V)\arrow[dl,orange]\arrow[ddl,green]\\
			&\varinjlim\mathcal{F}_i(U)\arrow[r,red]\arrow[d,"\varphi_U",blue]&\varinjlim\mathcal{F}_i(V)\arrow[d,"\varphi_V"',shift right=.2ex,red]\arrow[d,shift left=.2ex,orange]&\\
			&\mathcal{G}(U)\arrow[r,shift left=.2ex,blue]\arrow[r,shift right=.2ex,cyan]&\mathcal{G}(V)&
		\end{tikzcd}
		\hspace{20pt}
		\begin{tikzcd}
			\mathcal{F}_i\arrow[d,shift right=.2ex,red]\arrow[d,shift left=.2ex,blue]&\\
			\varinjlim\mathcal{F}_i\arrow[r,blue]\arrow[d,"\varphi",red]&\left(\varinjlim\mathcal{F}_i\right)^+\arrow[ld,"\overline{\varphi}",blue]&\\
			\mathcal{G}&
		\end{tikzcd}
	\end{equation*}
	To check it is a sheaf, use \Cref{1.2} to obtain the right diagram.
\end{exercise}

\begin{exercise}\label{ex1.11}
	Note that \textcolor{cyan}{Noetherian spaces are quasi-compact}. For the locality part, Let $s,t\in\varinjlim\mathcal{F}_i(U)$ and $s|_{V_k}=t|_{V_k}$. By \textcolor{cyan}{Atiyah-Macdonald Exercise 2.15(i)}, we have $s_{i_k}|_{V_k}\in\mathcal{F}_{i_k}\!\left(V_k\right)$ and $t_{j_k}|_{V_k}\in\mathcal{F}_{j_k}\!\left(V_k\right)$. Find an index $l\geqslant i_k,j_k$ for all $k$ due to there are finitely many $V_k$, so $\mu_l\!\left(s_l|_{V_k}\right)=\mu_l\!\left(t_l|_{V_k}\right)$. By \textcolor{cyan}{Atiyah-Macdonald Exercise 2.15(ii)}, we can find $l^\prime\geqslant l$ such that $s_{l^\prime}|_{V_k}=t_{l^\prime}|_{V_k}$. Then use the locality of $\mathcal{F}_l(U)$. The gluing part is similar.
\end{exercise}

\begin{exercise}\label{ex1.12}
	Similar to \Cref{ex1.10}.
\end{exercise}

\begin{exercise}\label{ex1.13}
	\begin{enumerate}
		\item[$\Rightarrow$:] Let $V_t\subseteq E(\mathcal{F})$ be a basic open set, and let $x\in s^{-1}\!\left(V_t\right)$, then $s(x)\in V_t$. By hypothesis, there exist a neighborhood $W\subseteq U$ of $x$ and $r\in\mathcal{F}(W)$ such that $s(y)=r_y$ for all $y\in W$. Take $W^\prime=W\cap V$. For each $y\in W^\prime$, $r_y=t_y\in V_t$, so $W^\prime$ is an open subset of $s^{-1}\!\left(V_t\right)$.
		\item[$\Leftarrow$:] Since $s(x)\in\mathcal{F}_x$, by \textcolor{cyan}{Atiyah-Macdonald Exercise 2.15(i)}, $s(x)=t_x$ for some $t\in\mathcal{F}(V)$, so $V_t$ is an open neighborhood of $s(x)$, and therefore $s^{-1}\!\left(V_t\right)$ is an open neighborhood of $x$. For every $y\in s^{-1}\!\left(V_t\right)$, by definition we have $s(y)\in V_t$ and $s(y)\in\mathcal{F}_y$, so $s(y)=t_y$.
	\end{enumerate}
\end{exercise}

\begin{exercise}\label{ex1.14}
	$\mu_U(s)=s_x$. By \textcolor{cyan}{Atiyah-Macdonald Exercise 2.15(ii)}, $\mu_{UV}(s)=\mathcal{F}\!\left(\iota_{VU}\right)(s)=0$ for some $V$.
\end{exercise}

\begin{exercise}\label{ex1.15}
	Let $\left\{U_i\right\}$ be an open cover of $U$, and let $\varphi_i:\mathcal{F}|_{U_i}\rightarrow\mathcal{G}|_{U_i}$ such that $\varphi_i|_{U_i\cap U_j}=\varphi_j|_{U_i\cap U_j}$. For each $V\subseteq U$, we have $\mathcal{F}|_U(V)=\mathcal{F}(V)$. Let $s\in\mathcal{F}(V)$ and denote $V_i=V\cap U_i$, then $\varphi_i\!\left(s|_{V_i}\right)\in\mathcal{G}|_{U_i}(V_i)$ satisfy $\left.\varphi_i\!\left(s|_{V_i}\right)\middle|_{V_i\cap V_j}\right.=\left.\varphi_i\middle|_{V_i\cap V_j}\!\left(s|_{V_i\cap V_j}\right)\right.=\left.\varphi_j\!\left(s|_{V_j}\right)\middle|_{V_i\cap V_j}\right.$. Using the gluing property of $\mathcal{G}$ we have a morphism $\varphi_V:\mathcal{F}|_U(V)\rightarrow\mathcal{G}|_U(V)$. Thus there is $\varphi:\mathcal{F}|_U\rightarrow\mathcal{G}|_U$. The locality is similar.
\end{exercise}

\begin{exercise}\label{ex1.16}
	\begin{enumerate}[label=(\alph*),ref=\theexercise(\alph*)]\crefalias{enumi}{exercise}
		\addtocounter{enumi}{1}
		\item\label{ex1.16.2} Denote $0\rightarrow\mathcal{F}^\prime\xrightarrow{\varphi}\mathcal{F}\xrightarrow{\psi}\mathcal{F}^{\prime\prime}\rightarrow0$. Similar to \Cref{ex1.8}, we have $\psi_{V_x\cap V_y}\!\left(t_x|_{V_x\cap V_y}\right)=\psi_{V_x\cap V_y}\!\left(t_y|_{V_x\cap V_y}\right)$. Since $\operatorname{im}\!\left(\varphi_{V_x\cap V_y}\right)=\ker\!\left(\psi_{V_x\cap V_y}\right)$, there is $r_{xy}\in\mathcal{F}^\prime\!\left(V_x\cap V_y\right)$ such that $\varphi_{V_x\cap V_y}\!\left(r_{xy}\right)=t_x|_{V_x\cap V_y}-t_y|_{V_x\cap V_y}$. By the flasqueness of $\mathcal{F}^\prime$, $\varphi_{V_x}\!\left(r_x\right)|_{V_x\cap V_y}=\varphi_{V_x\cap V_y}\!\left(r_{xy}\right)$ for some $r_x\in\mathcal{F}^\prime\!\left(V_x\right)$. Take $t_x^\prime=t_x-\varphi_{V_x}\!\left(r_x\right)$, then $t_x^\prime$ and $t_y$ are compatible on $V_x\cap V_y$.
		\item\label{ex1.16.3} Use \hyperref[ex1.16.2]{(b)} and \textcolor{cyan}{five lemma}.
		\addtocounter{enumi}{1}
		\item\label{ex1.16.5} If $t\in\mathcal{G}(V)$, let $s|_V=t$ and $s=0$ at $U\backslash V$. Define $\mathcal{F}(U)\rightarrow\mathcal{G}(U)$ to be $s\mapsto\left(x\mapsto s_x\right)$.
	\end{enumerate}
\end{exercise}

\begin{exercise}\label{ex1.17}
	Prove that if $y\in\overline{\{x\}}$, then $y\in U$ implies $x\in U$. It is a sheaf due to $x$ has to simultaneously be in or not be in all the open sets which make up an open cover.
\end{exercise}

\begin{exercise}\label{ex1.18}
	Apply the universal properties of direct limit and of sheafification on the restriction map $\mathcal{F}\!\left(f^{-1}(V)\right)\rightarrow\mathcal{F}(U)$ for $V\supseteq f(U)$. We also have $\displaystyle G(V)\xrightarrow{\mu_V}\varinjlim_{V^\prime\supseteq U^\prime}G\!\left(V^\prime\right)\rightarrow\left(\varinjlim_{V^\prime\supseteq U^\prime}G\!\left(V^\prime\right)\right)^+$ for $V\supseteq U^\prime$, where $U^\prime=f\!\left(f^{-1}(U)\right)$.
\end{exercise}

\begin{exercise}\label{ex1.19}
	\begin{enumerate}[label=(\alph*),ref=\theexercise(\alph*)]\crefalias{enumi}{exercise}
		\item\label{ex1.19.1} $i^{-1}(V)=V\cap Z$, and $V\cap Z$ is open in the subspace $Z$.
		\item\label{ex1.19.2} Use \Cref{1.1} to prove uniqueness.
		\item\label{ex1.19.3} Work on stalks.
	\end{enumerate}
\end{exercise}

\begin{exercise}\label{ex1.20}
	\begin{enumerate}[label=(\alph*),ref=\theexercise(\alph*)]\crefalias{enumi}{exercise}
		\item\label{ex1.20.1} $\Gamma_{Z\cap U}(U,\mathcal{F}|_U)=\left\{s\in\mathcal{F}(U)\middle|s_x\neq0\text{ for all }x\in Z\cap U\right\}$. Be aware that we only get $s\in \mathcal{F}(U)$ from the gluing axiom, then use support to derive $s\in\Gamma_{Z\cap U}(U,\mathcal{F}|_U)$.
		\item\label{ex1.20.2} $\displaystyle j_*\!\left(\mathcal{F}|_U\right)(V)=\left(\varinjlim_{W\supseteq U\cap V}\mathcal{F}(W)\right)^+=\mathcal{F}(U\cap V)$.
	\end{enumerate}
\end{exercise}

\begin{exercise}\label{ex1.21}
	
\end{exercise}

\begin{exercise}\label{ex1.22}
	$\mathcal{F}(U)=\left\{\left(s_i\right)\in\prod_{i\in I}\mathcal{F}_i\!\left(U\cap U_i\right)\middle|\varphi_{ij}\!\left(s_i|_{U\cap U_i\cap U_j}\right)=s_j|_{U\cap U_i\cap U_j}\right\}$, $\psi_i:\left(s_j\right)_{j\in I}\mapsto s_i$.
\end{exercise}

	\section{Schemes}

\begin{exercise}\label{ex2.1}
	\textcolor{cyan}{Atiyah-Macdonald Exercise 3.21(i)}.
\end{exercise}

\begin{exercise}\label{ex2.2}
	Let $X=\bigcup_iV_i$ and each $\left(V_i,\mathcal{O}_X|_{V_i}\right)$ is an affine scheme $\operatorname{Spec}\!\left(A_i\right)$. Since $U=\bigcup_jU_{f_j}$ is the union of some basic open sets, use \Cref{ex2.1} so $\left(V_i\cap U_{f_j},\mathcal{O}_X|_{V_i\cap U_{f_j}}\right)$ is $\operatorname{Spec}\!\left(\left(A_i\right)_{f_j}\right)$.
\end{exercise}

\begin{exercise}\label{ex2.3}
	
\end{exercise}

\begin{exercise}\label{ex2.4}
	
\end{exercise}

\begin{exercise}\label{ex2.5}
	
\end{exercise}

\begin{exercise}\label{ex2.6}
	
\end{exercise}

\begin{exercise}\label{ex2.7}
	
\end{exercise}

\begin{exercise}\label{ex2.8}
	\begin{enumerate}
		\item[(ii)] Use \Cref{2.15}.
	\end{enumerate}
\end{exercise}

\begin{exercise}\label{ex2.9}
	
\end{exercise}

\begin{exercise}\label{ex2.10}
	
\end{exercise}

\begin{exercise}\label{ex2.11}
	\begin{enumerate}
		\item[(i),(ii)] Consider a residue field and pass to vector spaces. Use the right exactness of tensor product.
		\item[(iii)] If $m>n$, let $\psi:A^m\xrightarrow{\phi}A^n\hookrightarrow A^m$. By \Cref{2.4}, $\psi$ satisfies $\psi^n+a_{n-1}\psi^{n-1}+\cdots+a_0=0$. Pick such equation to have the minimum degree, then $\psi$ injective implies $a_0\neq0$. Apply the equation to element $(0,\cdots,0,1)$ to obtain contradiction.
	\end{enumerate}
\end{exercise}

\begin{exercise}\label{ex2.12}
	
\end{exercise}

\begin{exercise}\label{ex2.13}
	
\end{exercise}

\begin{exercise}\label{ex2.14}
	
\end{exercise}

\begin{exercise}\label{ex2.15}
	\begin{enumerate}[label=(\roman*),ref=\theexercise(\roman*)]\crefalias{enumi}{exercise}
		\item\label{ex2.15.1} Use \Cref{ex2.14} to increase indices.
		\item\label{ex2.15.2} \textcolor{magenta}{$x_i=\sum_k\left(y_k-\mu_{kl}\left(y_k\right)\right)$ is an equation in $\bigoplus M_i$, so $\mu_{ij}$ cannot be applied directly.} Use canonical projection $\pi_i:\bigoplus M_i\rightarrow M_i$ on the equation first.
	\end{enumerate}
\end{exercise}

\begin{exercise}[Universal property of direct limit]\label{ex2.16}
	
\end{exercise}

\begin{exercise}\label{ex2.17}
	
\end{exercise}

\begin{exercise}\label{ex2.18}
	
\end{exercise}

\begin{exercise}\label{ex2.19}
	
\end{exercise}

	\section{First Properties of Schemes}

\begin{exercise}\label{ex3.1}
	Let $V_i=\operatorname{Spec}\!\left(B_i\right)$ be an affine open cover, and let $f^{-1}\!\left(V_i\right)$ be covered by $U_{ij}=\operatorname{Spec}\!\left(A_{ij}\right)$, where $A_{ij}$ is a finitely generated $B_i$-algebra. Let $V_i\cap V=\bigcup_kD_V\!\left(b_{ik}\right)$ where $b_{ik}\in B$, then $D_V\!\left(b_{ik}\right)\hookrightarrow V_i$ induces a map $B_i\rightarrow B_{b_{ik}}$. Also, $U_{ij}\cap f^{-1}\!\left(D_V\!\left(b_{ik}\right)\right)=D_{U_{ij}}\!\left(f_{U_{ij}}^\#\left(b_{ik}\right)\right)$. Thus, $f^{-1}(V)$ can be covered by $\operatorname{Spec}\!\left(\left(A_{ij}\right)_{f_{U_{ij}}^\#\!\left(b_{ik}\right)}\right)$, where $\left(A_{ij}\right)_{f_{U_{ij}}^\#\!\left(b_{ik}\right)}=A_{ij}\otimes_{B_i}B_{b_{ik}}$ is a finitely generated $B_{b_{ik}}$-algebra.
\end{exercise}

\begin{exercise}\label{ex3.2}
	Use the expression $U_{ij}\cap f^{-1}\!\left(D_V\!\left(b_{ik}\right)\right)$ in \Cref{ex3.1}. Since $V=\bigcup_{i,k}D_V\!\left(b_{ik}\right)$ is quasi-compact, so $i$, $k$ can be chosen to have finitely many. $j$ has finitely many by hypothesis.
\end{exercise}

\begin{exercise}\label{ex3.3}
	\begin{enumerate}[label=(\alph*),ref=\theexercise(\alph*)]\crefalias{enumi}{exercise}
		\addtocounter{enumi}{2}
		\item\label{ex3.3.3} Let $U_i=\operatorname{Spec}\!\left(A_i\right)$ be an affine open cover of $f^{-1}(V)$ such that $A_i$ is a finitely generated $B$-algebra. Let $U_i\cap U=\bigcup_kD_U\!\left(a_{ik}\right)$, where $a_{ik}\in A$. Since $U$ is affine, the union $U=\bigcup_{i,k}D_U\!\left(a_{ik}\right)$ can be finite. Since each $D_U\!\left(a_{ik}\right)$ is affine, there are finitely many $a_{ikl}^\prime\in A_i$ such that $D_U\!\left(a_{ik}\right)=\bigcup_lD_{U_i}\!\left(a_{ikl}^\prime\right)$.
		
		The disjoint union of all inclusions $D_{U_i}\!\left(a_{ikl}^\prime\right)\subseteq D_U\!\left(a_{ik}\right)$ induces an injective map
		\begin{equation*}
			\varphi:A_{a_{ik}}\rightarrow\prod_l\left(A_i\right)_{a_{ikl}^\prime}\qquad \alpha\mapsto\left(\alpha|_{D_{U_i}\!\left(a_{ikl}^\prime\right)}\right)_l.
		\end{equation*}		
		However, if $\left(\beta_l\right)\in\prod_l\left(A_i\right)_{a_{ikl}^\prime}$ is an image from $A_{a_{ik}}$, the sheaf property on $D_U\!\left(a_{ik}\right)$ implies $\beta_{l_1}|_{D_{U_i}\!\left(a_{ikl_1}^\prime a_{ikl_2}^\prime\right)}=\beta_{l_2}|_{D_{U_i}\!\left(a_{ikl_1}^\prime a_{ikl_2}^\prime\right)}$, so $A_{a_{ik}}$ is isomorphic to a quotient ring of $\prod_l\left(A_i\right)_{a_{ikl}^\prime}$. The ring $\prod_l\left(A_i\right)_{a_{ikl}^\prime}$ is a finitely generated $B$-algebra, so is $A_{a_{ik}}$. By a similar method, $A$ is a finitely generated $B$-algebra.
	\end{enumerate}
\end{exercise}

\begin{exercise}\label{ex3.4}
	Similar to \Cref{ex3.1}.
\end{exercise}

\begin{exercise}\label{ex3.5}
	\begin{enumerate}[label=(\alph*),ref=\theexercise(\alph*)]\crefalias{enumi}{exercise}
		\item \textcolor{cyan}{Atiyah-Macdonald Exercise 8.4}.
		\item Let $V_i=\operatorname{Spec}\!\left(B_i\right)$, $f^{-1}\!\left(V_i\right)=\operatorname{Spec}\!\left(A_i\right)$. $B_i\rightarrow A_i$ is an integral homomorphism. If $C\subseteq X$ is closed, by \textcolor{cyan}{Atiyah-Macdonald Exercise 5.1}, $f\!\left(C\cap f^{-1}\!\left(V_i\right)\right)=f(C)\cap V_i$ is closed in subspace $V_i$, so $f(C)$ is closed.
		\item $k[x]\rightarrow k\!\left[x,x^{-1}\right]\times k$ defined to be $k[x]\hookrightarrow k\!\left[x,x^{-1}\right]$ and $x\mapsto0$ in $k[x]\rightarrow k$.
	\end{enumerate}
\end{exercise}

\begin{exercise}\label{ex3.6}
	By \Cref{3.1} and \Cref{ex2.9}, $\xi$ is unique. Also, $\xi$ is contained in every open subset. For the $U=\operatorname{Spec}A$, $A=\mathcal{O}_X(U)$ is an integral domain, so $\mathcal{O}_{X,\xi}=\left(\mathcal{O}_X|_U\right)_{\xi}=A_{(0)}$ is a field.
\end{exercise}

\begin{exercise}\label{ex3.7}
	Let $V=\operatorname{Spec}B\subseteq Y$ and $U_i=\operatorname{Spec}\!\left(A_i\right)\subseteq f^{-1}(V)$, such that $A_i$ is a finitely generated $B$-algebra. Since $X$ is irreducible, $U_i$ is dense in $X$, so $f\!\left(U_i\right)$ is dense in $V$, and by \textcolor{cyan}{Atiyah-Macdonald Exercise 1.21(v)}, $B\rightarrow A_i$ is injective.
	
	Let $S=B-\{0\}$. Then $S^{-1}A_i$ is a finitely generated $k(B)$-algebra. By \textcolor{cyan}{Noether's normalization lemma}, $S^{-1}A_i$ is integral over $k(B)\!\left[t_1,\dots,t_n\right]$. By \textcolor{cyan}{Atiyah-Macdonald Exercise 5.10}, $\operatorname{Spec}\!\left(S^{-1}A_i\right)\rightarrow\operatorname{Spec}\!\left(k(B)\!\left[t_1,\dots,t_n\right]\right)$ is surjective. Since $f$ is generically finite, $\operatorname{Spec}\!\left(S^{-1}A_i\right)=\operatorname{Spec}\!\left(A_i\otimes_Bk(B)\right)$ is a finite set, so $n=0$. Therefore, $S^{-1}A_i$ is a finite dimensional $k(B)$-vector space, so $S^{-1}A_i=k(A)$, which is a finite field extension of $k(B)$.
	
	Let $x_{ij}$ generate $A_i$ over $B$. Then $x_{ij}$ satisfies an algebraic dependence equation whose coefficients are all in $B$. Let $b_{ij}$ be the leading coefficient and let $s_i$ be their multiplication with respect to $j$. Then each $x_{ij}\in A_i\!\left[s_i^{-1}\right]$ are integral, so $A_i\!\left[s_i^{-1}\right]$ is a finitely generated $B\!\left[s_i^{-1}\right]$-module.
	
	Since $X$ is irreducible, $W=\bigcap_iU_i$ is a non-empty open set. There exists $a_i\in A_i\!\left[s_i^{-1}\right]$ such that $D\!\left(a_i\right)\in W$. Since $A_i\!\left[s_i^{-1}\right]$ is finite over $B\!\left[s_i^{-1}\right]$, take $b_i$ to be the constant coefficient of the monic equation, so $b_i\in\left(a_i\right)$, and therefore $D\!\left(b_i\right)\subseteq D\left(a_i\right)\subseteq W$. Let $b$ be the multiplication of $b_i$ with respect to $i$, then $U=D(b)$ is the desired open set.
\end{exercise}

\begin{exercise}\label{ex3.8}
	Let $V=\operatorname{Spec}B$ and let $x\in U\cap V$. Then there are $f\in A$ and $g\in B$ such that $D_V(g)\subseteq D_U(f)\subseteq U\cap V$ with $x\notin D_U(f)$ and $x\notin D_V(g)$. The restriction $D_U(f)\subseteq V$ induces a homomorphism $B\rightarrow A_f$ in which $g\mapsto a/f^n$. By \Cref{ex2.16.1}, $D_U(f)\cap D_V(g)=D_{\operatorname{Spec}\left(A_f\right)}\!\left(a/f^n\right)=D_U(af)$, so $D_V(g)=D_U(af)$. Therefore, $U\cap V$ can be covered by $\left\{W_i\right\}$ such that $W_i=D_U\!\left(f_i\right)=D_V\!\left(g_i\right)$. Since \textcolor{cyan}{localization commutes with integral closure}, the isomorphisms $\varphi_i:D_{\tilde{U}}\!\left(f_i\right)\rightarrow D_{\tilde{V}}\!\left(g_i\right)$ are compatible on overlap. Then use \Cref{ex2.12}.
	
	Break to affine case, we need to prove that the integral closure of a finitely generated $k$-algebra $A$ is a finitely generated $A$-module. By \textcolor{cyan}{Noether's normalization lemma}, $A$ is finite over $k\!\left[y_1,\dots,y_r\right]$. Let $K=\operatorname{Frac}\!\left(k\!\left[y_1,\dots,y_r\right]\right)$ and $L=\operatorname{Frac}A$. Then $L/K$ is a finite field extension. Let $B$ be the integral closure of $k\!\left[y_1,\dots,y_r\right]$ in $L$. \textcolor{red}{Generalization of} \textcolor{cyan}{Atiyah-Macdonald Exercise 5.15}.
\end{exercise}

\begin{exercise}\label{ex3.9}
	Use \Cref{3.3}. $k(s)\otimes_kk(t)$ is a PID. The prime ideals are $(0)$ and ideals generate by the elements which is irreducible and cannot be written in the form $f(s)\otimes_kg(t)$.
\end{exercise}

\begin{exercise}\label{ex3.10}
	\begin{enumerate}[label=(\alph*),ref=\theexercise(\alph*)]\crefalias{enumi}{exercise}
		\item\label{ex3.10.1} \textcolor{cyan}{Atiyah-Macdonald Exercise 3.21(iv)} is the affine case. Let $\left\{U_i\right\}$ be an affine open cover of $X$, then each $\varphi_i:\left(U_i\right)_y\rightarrow U_i\cap f^{-1}(y)$ is a homeomorphism. By the universal property of fiber product, the morphisms in the diagram
		\begin{equation*}
			\begin{tikzcd}
				\left(U_i\cap U_j\right)_y\arrow[r]\arrow[d]\arrow[rd]&\left(U_i\right)_y\arrow[d]\\
				\left(U_j\right)_y\arrow[r]&X_y
			\end{tikzcd}
		\end{equation*}
		are all unique, so the diagram commutes, and therefore $\varphi_i$ and $\varphi_j$ are compatible on overlap.
		\item\label{ex3.10.2} $X_y=\operatorname{Spec}\!\left(k[t]/\!\left(t^2-a\right)\right)$. $X_\eta=\operatorname{Spec}\!\left(k(s)[t]/\!\left(t^2-s\right)\right)$. $k$ is algebraically closed, so $t^2-s$ is irreducible.
	\end{enumerate}
\end{exercise}

\begin{exercise}\label{ex3.11}
	\begin{enumerate}[label=(\alph*),ref=\theexercise(\alph*)]\crefalias{enumi}{exercise}
		\item\label{ex3.11.1} Use \textcolor{cyan}{Atiyah-Macdonald Exercise 2.2} for the affine case. When only $X$ is affine, use a method similar to \Cref{ex3.10.1} to glue morphisms. Let $g:X^\prime\rightarrow X$, $X_i$ an affine cover of $X$, $Y_i=f^{-1}\!\left(X_i\right)$, and $X^\prime_i=g^{-1}\!\left(X_i\right)$. From the commutative diagram of $Y\times_XX^\prime_i$, the image of $Y\times_XX^\prime_i\rightarrow Y$ is inside $Y_i$. By the uniqueness from \Cref{3.3}, we have $Y\times_XX^\prime_i\cong Y_i\times_XX^\prime_i\cong Y_i\times_{X_i}X^\prime_i$. Then glue $X^\prime_i$.
		\item\label{ex3.11.2} Let $D(f)$ be a basic open set of $X$ and let $y\in Y\cap D(f)$, then $f\notin\mathfrak{p}_{i(y)}$. The map $i_y^\#:\mathcal{O}_{X,i(y)}\rightarrow\mathcal{O}_{Y,y}$ is a surjective local homomorphism, which implies $\left(i^\#(f)\right)_y=i_y^\#\!\left(f_{i(y)}\right)\notin\mathfrak{m}_{Y,y}$. By \Cref{ex2.16.1}, $U\cap Y\cap D(f)$ is a basic open set of $Y$ for each open affine $U\subseteq Y$. Since a closed subset of a quasi-compact space is quasi-compact, the cover of $Y$ can use only finitely many $f_i$. The set $\left\{f_i\right\}$ can be expanded so it generates the unit ideal of $A$. Restrict each $f_i$ on $Y$ and use \Cref{ex2.17.2}.
		\item\label{ex3.11.3} Let $U_i=\operatorname{Spec}\!\left(A_i\right)$ be an affine cover of $X$. By \hyperref[ex3.11.2]{(b)}, $Y^\prime\cap U_i=V\!\left(\mathfrak{a}_i\right)$; by hypothesis, $Y\cap U_i=V\!\left(\!\sqrt{\mathfrak{a}_i}\right)$. There is a natural homomorphism $A/\mathfrak{a}_i\rightarrow\left.A\middle/\!\sqrt{\mathfrak{a}_i}\right.$, which induces a morphism $\varphi_i:Y\cap U_i\rightarrow Y^\prime\cap U_i$. Therefore, we have a morphism of sheaves $\mathcal{O}_{U_i}|_{Y^\prime\cap U_i}\rightarrow\mathcal{O}_{U_i}|_{Y\cap U_i}$, which is exactly $\mathcal{O}_X|_{Y^\prime\cap U_i}\rightarrow\mathcal{O}_X|_{Y\cap U_i}$. Restrict on $U_i\cap U_j$, so $\varphi_i$ and $\varphi_j$ are compatible on overlap.
		\item\label{ex3.11.4} Let $\mathfrak{a}_i=\ker\!\left(A_i\rightarrow\Gamma\!\left(f^{-1}\!\left(U_i\right),\mathcal{O}_Z\right)\right)$, and let $Y_i=\operatorname{Spec}\!\left(A_i/\mathfrak{a}_i\right)$. By \Cref{ex2.4}, the ring homomorphism $A_i/\mathfrak{a}_i\rightarrow\Gamma\!\left(f^{-1}\!\left(U_i\right),\mathcal{O}_Z\right)$ induces a morphism $\varphi_i:f^{-1}\!\left(U_i\right)\rightarrow Y_i$. Combine $\varphi_i$ with $Y_i\rightarrow U_i$ and use \Cref{ex2.4} again to prove the corresponding diagram commutes. Combine $\varphi_i$ with $Y_i\rightarrow Y$ and use a method similar to \hyperref[ex3.11.3]{(c)}, $\varphi_i$ can be glued.
	\end{enumerate}
\end{exercise}

\begin{exercise}\label{ex3.12}
	\begin{enumerate}[label=(\alph*),ref=\theexercise(\alph*)]\crefalias{enumi}{exercise}
		\item Similar to \textcolor{cyan}{Atiyah-Macdonald Exercise 1.21(iv)}, use isomorphism $S/\ker\varphi\rightarrow T$.
		\item Use \Cref{ex2.14.3}.
	\end{enumerate}
\end{exercise}

\begin{exercise}\label{ex3.13}
	\begin{enumerate}[label=(\alph*),ref=\theexercise(\alph*)]\crefalias{enumi}{exercise}
		\item 
		\item 
		\item 
		\item 
		\item 
		\item 
		\item 
	\end{enumerate}
\end{exercise}

\begin{exercise}\label{ex3.14}
	\textcolor{cyan}{Atiyah-Macdonald Exercise 5.24(ii) and 5.26}. A counter example is $\mathbb{Z}$.
\end{exercise}

\begin{exercise}\label{ex3.15}
	\begin{enumerate}[label=(\alph*),ref=\theexercise(\alph*)]\crefalias{enumi}{exercise}
		\item 
		\item 
		\item 
	\end{enumerate}
\end{exercise}

\begin{exercise}\label{ex3.16}
	Suppose there is a closed subset such that $\mathcal{P}$ doesn't hold. Let $\Sigma$ be the collection of closed subsets of $X$ such that $\mathcal{P}$ doesn't hold, then $\Sigma$ is not empty. Since $X$ is Noetherian, $\Sigma$ has a minimal element $Z$. However, $\mathcal{P}$ holds for each closed subset of $Z$, which contradicts to the hypothesis.
\end{exercise}

\begin{exercise}\label{ex3.17}
	\begin{enumerate}[label=(\alph*),ref=\theexercise(\alph*)]\crefalias{enumi}{exercise}
		\item \Cref{ex2.9}.
		\item If $C\subseteq X$ is minimal closed, the closure of any point in $C$ is $C$ itself.
		\item The closure of a point $x$ is irreducible with $x$ the generic point. Prove that $x\in\overline{\{y\}}$ and $y\in\overline{\{x\}}$ simultaneously is impossible.
		\addtocounter{enumi}{2}
		\item 
	\end{enumerate}
\end{exercise}

\begin{exercise}\label{ex3.18}
	\begin{enumerate}[label=(\alph*),ref=\theexercise(\alph*)]\crefalias{enumi}{exercise}
		\item 
		\item 
		\item 
		\item 
	\end{enumerate}
\end{exercise}

\begin{exercise}\label{ex3.19}
	\begin{enumerate}[label=(\alph*),ref=\theexercise(\alph*)]\crefalias{enumi}{exercise}
		\item 
		\item 
		\item 
	\end{enumerate}
\end{exercise}

\begin{exercise}\label{ex3.20}
	\begin{enumerate}[label=(\alph*),ref=\theexercise(\alph*)]\crefalias{enumi}{exercise}
		\item 
		\item 
		\item 
		\item 
		\item 
		\item 
	\end{enumerate}
\end{exercise}

\begin{exercise}\label{ex3.21}
	
\end{exercise}

\begin{exercise}\label{ex3.22}
	\begin{enumerate}[label=(\alph*),ref=\theexercise(\alph*)]\crefalias{enumi}{exercise}
		\item 
		\item 
		\item 
		\item 
		\item 
	\end{enumerate}
\end{exercise}

\begin{exercise}\label{ex3.23}
	
\end{exercise}

	\section{Separated and Proper Morphisms}

\begin{exercise}\label{ex4.1}
	Use \textcolor{cyan}{Atiyah-Macdonald Corollary 2.17 and Exercise 5.1} and \hyperref[affinebaseext]{Affine base extension criterion} for universally closed. The global section of $\left(\Delta|_{U_i}\right)^\#$ is defined to be $a\otimes a^\prime\mapsto aa^\prime$, which is surjective. By \textcolor{cyan}{Atiyah-Macdonald Exercise 1.21(iv)}, $\Delta|_{U_i}$ is a closed immersion, so is $\Delta$.
\end{exercise}

\begin{exercise}\label{ex4.2}
	
\end{exercise}

\begin{exercise}\label{ex4.3}
	For any element $x\in A$, $x(1-ax)=0\in\mathfrak{q}$.
\end{exercise}

\begin{exercise}\label{ex4.4}
	$\mathfrak{m}=\sqrt{\mathfrak{q}}\supsetneq\mathfrak{q}\supsetneq\mathfrak{m}^2$.
\end{exercise}

\begin{exercise}\label{ex4.5}
	$\mathfrak{m}^2$ is $\mathfrak{m}$-primary.
\end{exercise}

\begin{exercise}\label{ex4.6}
	Use Urysohn's lemma to proof that if $f\left(x_1\right)=f\left(x_2\right)=0$ for all $f\in\mathfrak{q}$, $\mathfrak{q}$ cannot be primary, so there are infinite many minimal prime ideals. Now use \Cref{1.11.2} to proof that $(0)=\sqrt{(0)}=\bigcap\mathfrak{p}$ is the intersection of ALL minimal prime ideals.
\end{exercise}

\begin{exercise}\label{ex4.7}
	\begin{enumerate}
		\item[(ii),(iv),(v)] $(\mathfrak{a}[x])^c=\mathfrak{a}$.
		\item[(iii)] Let $\overline{f},\overline{g}\in\left(A/\mathfrak{q}\right)[x]$. If $\overline{f}\overline{g}=0$ but $\overline{g}\neq0$, by \Cref{ex1.2.3}, $\overline{a}\overline{f}=0$, so $aa_i\in\mathfrak{q}$. Hence $a_i\in\mathfrak{p}$ and $a_i^{n_i}\in\mathfrak{q}$, so $\overline{a_i}^{n_i}=0$ and use \Cref{ex1.2.2}.
	\end{enumerate}
\end{exercise}

\begin{exercise}\label{ex4.8}
	$k\left[x_1,\cdots,x_n\right]/\mathfrak{p_i}=k\left[x_{i+1},\cdots,x_n\right]$ integer domain, so $\mathfrak{p}_i$ prime. Note that $\mathfrak{p}_i^n$ is the set of all homogeneous polynomials in $k\left[x_1,\cdots,x_i\right]$ of degree $\geqslant n$.
\end{exercise}

\begin{exercise}\label{ex4.9}
	\begin{enumerate}[label=(\roman*),ref=\theexercise(\roman*)]\crefalias{enumi}{exercise}
		\item $\Leftarrow$: $x$ is contained in all prime ideals containing $(0:a)$, so $x\in\sqrt{(0:a)}$.
		\item If $\mathfrak{a}\subseteq\mathfrak{p}$ but $\mathfrak{a}^e\subseteq S^{-1}\mathfrak{p}^\prime\subsetneq S^{-1}\mathfrak{p}$, then $\mathfrak{a}\subseteq\mathfrak{a}^{ec}\subseteq\mathfrak{p}^\prime\subsetneq\mathfrak{p}$. If $\mathfrak{b}\subseteq S^{-1}\mathfrak{p}$ but $\mathfrak{b}^c\subseteq\mathfrak{p}^\prime\subsetneq\mathfrak{p}$, use $\mathfrak{b}^{ce}=\mathfrak{b}$ derived from \Cref{3.11.1}. $S^{-1}\mathfrak{p}$ is a prime ideal if $S$ doesn't meet $\mathfrak{p}$.
		\item $(0:a)=\bigcap\left(\mathfrak{q}_i:a\right)\subseteq\mathfrak{p}$. By \Cref{1.11.2}, $\left(\mathfrak{q}_i:a\right)\subseteq\mathfrak{p}$ for some $i$. Then use \Cref{4.4} and the minimality of $\mathfrak{p}$.
	\end{enumerate}
\end{exercise}

\begin{exercise}\label{ex4.10}
	\begin{enumerate}
		\item[(iv)] Let $x\neq0$, $(0:x)\subseteq\mathfrak{p}$ for some $\mathfrak{p}$. If $x\in S_\mathfrak{p}(0)$, then $sx=0$, so $s\in(0:x)$.
	\end{enumerate}
\end{exercise}

\begin{exercise}\label{ex4.11}
	Note that $x\in S(\mathfrak{a})$ $\Leftrightarrow$ there exists an $s\in S$ such that $sx\in\mathfrak{a}$. Show that $S_\mathfrak{p}(0)$ is contained in all $\mathfrak{p}$-primary ideals.\par
	If $0=\bigcap\mathfrak{q}_i$ is irredundant, $\sqrt{(0)}$ is the intersection of minimal prime ideals.
\end{exercise}

\begin{exercise}\label{ex4.12}
	Use \Cref{3.11.2}. Use \Cref{4.9} to show there are at most $2^n$ different $S(\mathfrak{a})$.
\end{exercise}

	\section{Integral Dependence and Valuations}

\begin{exercise}\label{ex5.1}
	$f^*(V(I))=\left\{f^{-1}(\mathfrak{q}):\mathfrak{q}\supseteq I\right\}$. Use \Cref{5.10} to proof $f^*(V(I))\supseteq V\left(f^{-1}(I)\right)$.
\end{exercise}

\begin{exercise}\label{ex5.2}
	$\mathfrak{p}=\ker f$ is prime. By \Cref{5.10}, there is a prime ideal $\mathfrak{q}\subseteq B$. $\operatorname{Frac}(B/\mathfrak{q})$ is an algebraic extension of $\operatorname{Frac}(A/\mathfrak{p})$, so $\operatorname{Frac}(A/\mathfrak{p})\rightarrow\Omega$ can be extended to $\operatorname{Frac}(B/\mathfrak{q})\rightarrow\Omega$.
\end{exercise}

\begin{exercise}\label{ex5.3}
	Use the operator $-\otimes_A c^n$.
\end{exercise}

\begin{exercise}\label{ex5.4}
	If $\left(\frac{1}{x+1}\right)^n+\sum\frac{f_i\left(x^2-1\right)}{g_i\left(x^2-1\right)(x+1)^i}=0$, multiply both side by $(x+1)^n$ and take $x=-1$ to get contradiction.
\end{exercise}

\begin{exercise}\label{ex5.5}
	\begin{enumerate}
		\item[(ii)] Use \Cref{5.10} to proof that the contraction of maximal ideals in $B$ are maximal in $A$.
	\end{enumerate}
\end{exercise}

\begin{exercise}\label{ex5.6}
	Let $b_i\in B_i$, $f_i\in A\left[x_i\right]$ monic, $f_i(b_i)=0$. Let $f=\left(f_1,\cdots,f_n\right)$, then $f(b_1,\cdots,b_n)=0$.
\end{exercise}

\begin{exercise}\label{ex5.7}
	Let $x\in B$ and $x$ integral over $A$. Let $x^n+a_{n-1}x^{n-1}+\cdots+a_0=0$ is the equation of the least degree. Analyze $x\left(x^{n-1}+a_{n-2}x^{n-2}+\cdots+a_1\right)=-a_0$.
\end{exercise}

\begin{exercise}\label{ex5.8}
	Proof that there is a ring $D\supseteq B$ such that $f$ and $g$ factor into linear factors in $D$. Use Kronecker's construction.
\end{exercise}

\begin{exercise}\label{ex5.9}
	Let $f^n+g_{n-1}f^{n-1}+\cdots+g_0=0$. Let $f_1=f-x^r$ with $r$ larger than all $\operatorname{deg}g_i$, then $f_1\left(f_1^{n-1}+h_{n-1}f_1^{n-2}+\cdots+h_1\right)=-h_0$ with $h_0$ monic. Then use \Cref{ex5.8}.
\end{exercise}

\begin{exercise}\label{ex5.10}
	\begin{enumerate}
		\item[(a) $\Rightarrow$ (b)] Let $f^*(V(\mathfrak{q}_1))=V(I)$, so $I\subseteq\mathfrak{p}_1\subseteq\mathfrak{p}_2$, and we have $\mathfrak{p}_2\in V(I)=f^*(V(\mathfrak{q}_1))$.
		\item[(a') $\Rightarrow$ (c')] By \hyperref[ex3.23]{\Cref*{ex3.23}(v)} and \Cref{ex3.26}, $f^*\left(\operatorname{Spec}\left(B_\mathfrak{q}\right)\right)=f^*\left(\varinjlim\operatorname{Spec}\left(B_t\right)\right)=\bigcap_tf_t^*\left(\operatorname{Spec}\left(B_t\right)\right)$. Each $f_t^*\left(\operatorname{Spec}\left(B_t\right)\right)$ is an open neighborhood of $\mathfrak{p}$, so it contains $A_\mathfrak{p}$.
	\end{enumerate}
\end{exercise}

\begin{exercise}\label{ex5.11}
	Use \Cref{ex3.18} and \Cref{ex5.10}.
\end{exercise}

\begin{exercise}\label{ex5.12}
	Note that $G$ has identity element so $x$ is a root of the polynomial $\prod_{\sigma\in G}\left(t-\sigma(x)\right)$. The coefficients are symmetric polynomials of $\sigma_i$. \textcolor{magenta}{$G$ acting on $S^{-1}A$ is defined to be $\sigma(x/s)=\sigma(x)/\sigma(s)$}.
\end{exercise}

\begin{exercise}\label{ex5.13}
	Let $\mathfrak{q}\in P$. Then $\sigma(\mathfrak{q})\cap A^G=\mathfrak{q}\cap A^G=\mathfrak{p}$, so $\sigma(\mathfrak{q})\in P$. Let $\mathfrak{q}_1$, $\mathfrak{q}_2\in P$ and $x\in\mathfrak{q}_1$. Then $\prod_{\sigma\in G}\sigma(x)\in\mathfrak{q}_1\cap A^G=\mathfrak{p}\subseteq\mathfrak{q}_2$, so $x\in\sigma^{-1}\left(\mathfrak{q}_2\right)$, and so $\mathfrak{q}_1\subseteq\bigcup_{\sigma\in G}\sigma\left(\mathfrak{q}_2\right)$. Use \Cref{1.11.1} and \Cref{5.9}.
\end{exercise}

\begin{exercise}\label{ex5.14}
	Note that $\sigma\in G$ fixes $A$.
\end{exercise}

\begin{exercise}\label{ex5.15}
	
\end{exercise}

\begin{exercise}\label{ex5.16}
	We have a homomorphism $f:k\left[y_1,\cdots,y_r\right]\hookrightarrow k\left[x_1,\cdots,x_n\right]\twoheadrightarrow  A$. Let $a=\left(a_1,\cdots,a_r\right)\in L$ corresponds to a maximal ideal $\mathfrak{m}_a=\left(y_1-a_1,\cdots,y_r-a_r\right)$. By \Cref{5.10}, we can find a prime ideal $\mathfrak{p}_a\subset A$ such that $f^{-1}(\mathfrak{p}_a)=\mathfrak{m}_a$, which is contained in a maximal $\mathfrak{n}_a$. Notice that $\mathfrak{n}_a$ correspond a point in $X$.
\end{exercise}

\begin{exercise}\label{ex5.17}
	Consider $k\left[t_1,\cdots,t_n\right]/I(X)$. Use \Cref{ex5.16}.
\end{exercise}

\begin{exercise}[\textcolor{red}{Zariski's lemma}]\label{ex5.18}
	Let $n$ be the number of generators. If $n=1$, $x_1$ is invertible, so $x_1^{-1}=\sum_{i=1}^nc_ix_1^i$.
	
	If $n>1$, let $A=k\left[x_1\right]$ and $K_A=\operatorname{Frac}(A)$. Then $K_A[x_2,\cdots,x_n]$ is a finite algebraic extension of $K_A$, so there are monic polynomials $f_i\in K_A[y]$ such that $f_i\left(x_i\right)=0$ for $2\leqslant i\leqslant n$. Pick $a\in A$ such that $af_2\cdots f_n\in A[y]$, then $x_2,\cdots,x_n$ are integral over $A_a$. Therefore, $K_A\subseteq B$ are integral over $A_a$.
	
	If $x_1$ transcendental over $k$, $A$ is a UFD. Since every UFD is integrally closed, $A$ is integrally closed, so $A_a$ is integrally closed by \Cref{5.12}. $K_A$ integral over $A_a$, so $A_a=K_A$, but such $a$ must be a unit in $A$.
\end{exercise}

\begin{exercise}\label{ex5.19}
	By \Cref{ex5.18}, $k\left[t_1,\cdots,t_n\right]/\mathfrak{m}$ is a finite algebraic extension of $k$, so $k\left[t_1,\cdots,t_n\right]/\mathfrak{m}\cong k$. Let $t_i\mapsto a_i$, then $t_i-a_i\in\mathfrak{m}$. Proof that $\left(t_1-a_1,\cdots,t_n-a_n\right)$ is maximal.
\end{exercise}

\begin{exercise}\label{ex5.20}
	Let $S=A-\{0\}$. By \hyperref[noethernormalization]{Noether normalization lemma}, $S^{-1}B$ is integral over $S^{-1}A\left[y_1,\cdots,y_r\right]$ with $y_1,\cdots,y_r\in S^{-1}B$. The denominators of $y_1,\cdots,y_r$ can be absorbed to $S^{-1}A$, so we can assume $y_1,\cdots,y_r\in B$. Let $B=A\left[x_1,\cdots,x_n\right]$, then $x_i/1\in S^{-1}B$ is integral over $S^{-1}A\left[y_1,\cdots,y_r\right]$. Reduce the coefficients of all integral dependence equations to a common denominator $s$, we have $x_i/1$ is integral over $A_s\left[y_1,\cdots,y_r\right]$. Therefore, $B_s=A_s\left[x_1/1,\cdots,x_n/1\right]$ is integral over $A_s\left[y_1,\cdots,y_r\right]$.
\end{exercise}

\begin{exercise}\label{ex5.21}
	We have
	\begin{equation*}
		\begin{tikzcd}
			A\arrow[r,"\phi_A"]\arrow[drrr,"f"{description,pos=0.35}]&A_s\arrow[r]\arrow[drr,"f_s"{description,pos=0.35}]&A_s\left[y_1,\cdots,y_n\right]\arrow[r]\arrow[dr,"f_s^\prime"{description,pos=0.4}]
			&B_s\arrow[d,"g"]&B\arrow[l,"\phi_B"']\\
			&&&\Omega&
		\end{tikzcd}
	\end{equation*}
	By \Cref{ex5.20}, $B_s$ integral over $A_s\left[y_1,\cdots,y_n\right]$, so $g$ is obtained from \Cref{ex5.2}. Then, $f_s$ by universal property of localization; $f_s^\prime$ by defining $y_i\mapsto0$. 
\end{exercise}

\begin{exercise}\label{ex5.22}
	Let $v\neq0$ an element of $B$. Then $B_v$ is integral over $A$, so $f:A\rightarrow\Omega$ can be extended to $\overline{f}:B_v\rightarrow\Omega$ by \Cref{ex5.21}, with $s\neq0$ an element of $A$ such that $f(s)\neq0$. Find an $\mathfrak{m}$ such that $s\notin\mathfrak{m}$ and use \Cref{ex5.21} again, $g:A\rightarrow\overline{A/\mathfrak{m}}$ can be extended to $\overline{g}:B_v\rightarrow\overline{A/\mathfrak{m}}$. Since $\overline{g}(s)\neq0$, $\overline{g}$ is not trivial, so $\overline{g}(1/v)\overline{g}(v)=\overline{g}(1)=1$ means $\overline{g}(v)\neq0$.
	
	Since $\overline{A/\mathfrak{m}}$ integral over $A/\mathfrak{m}$, $\overline{g}\left(B_v\right)$ integral over $A/\mathfrak{m}$, so $\overline{g}\left(B_v\right)\cong B_v/\ker\left(\overline{g}\right)$ is a field and therefore $\ker\overline{g}$ is maximal. By \Cref{5.10}, $\ker\overline{g}\cap B$ is maximal in $B$ that doesn't contain $v$.
\end{exercise}

\begin{exercise}\label{ex5.23}
	\begin{enumerate}
		\item[(i) $\Rightarrow$ (ii)] Note that $f:A\rightarrow A/\ker f$. Use \Cref{1.1}.
		\item[(ii) $\Rightarrow$ (iii)] Consider $A/\mathfrak{p}$.
		\item[(iii) $\Rightarrow$ (i)] If $\mathfrak{p}$ is not an intersection of maximal ideals, the Jacobson radical $\mathfrak{J}$ of $A/\mathfrak{p}$ is not zero. Pick $f\in\mathfrak{J}$, then $(A/\mathfrak{p})_f\neq0$. The contraction of a maximal ideal of $(A/\mathfrak{p})_f$ is $\mathfrak{q}\in A$ which is not maximal and doesn't contain $f$, but all prime ideals strictly containing $\mathfrak{q}$ contains $f$.
	\end{enumerate}
\end{exercise}

\begin{exercise}\label{ex5.24}
	\begin{enumerate}[label=(\roman*),ref=\theexercise(\roman*)]\crefalias{enumi}{exercise}
		\item Use \hyperref[ex5.23]{\Cref*{ex5.23}(iii)} and \hyperref[5.11]{going-up theorem}.
		\item Consider $B/\mathfrak{q}$ and $A/\mathfrak{q}^c$. Then use \Cref{ex5.22}, \hyperref[ex5.23]{\Cref*{ex5.23}(i)} and \Cref{1.1}.
	\end{enumerate}
\end{exercise}

\begin{exercise}\label{ex5.25}
	\begin{enumerate}
		\item[(i) $\Rightarrow$ (ii)] Similar to the proof of \Cref{ex5.22}, we can find $\overline{g}:B\rightarrow\overline{A/\mathfrak{m}}$. Let $B=A\left[x_1,\cdots,x_n\right]$, then $\overline{g}\left(x_i\right)$ integral over $A/\mathfrak{m}$, and therefore $\overline{g}(B)$ is finitely generated as an $A/\mathfrak{m}$-algebra. By \hyperref[ex5.18]{Zariski's lemma}, $\overline{g}(B)$ is finite over $A/\mathfrak{m}$. Note that $B$ is a field, so $\overline{g}$ injective.
		\item[(ii) $\Rightarrow$ (i)] Take $B=A/\mathfrak{p}$. For any $f\neq0$ in $B$, $B_f$ is a finitely generated $A$-algebra. If $B_f$ is a field, it is finite over $A$, so integral over $A$. The integral dependence equation also asserts that $B_f$ is integral over $B$, so $B$ is a field by \Cref{5.7}. If $B_f$ is not a field, it has a non-zero prime ideal whose contraction in $B$ doesn't contain $f$. Therefore, the intersection of all prime ideals in $B$ is zero.
	\end{enumerate}
\end{exercise}

\begin{exercise}\label{ex5.26}
	\begin{enumerate}
		\item[(2) $\Rightarrow$ (3)] $U_1\cap X_0=U_2\cap X_0$ $\Rightarrow$ $U_1^c\cap X_0=U_2^c\cap X_0$, then take closure.
		\item[(i) $\Rightarrow$ (ii) $\Rightarrow$ (iii)] Consider $V(\mathfrak{a})\cap X_f\cap\operatorname{Max}(A)$. If all maximal ideals in $V(\mathfrak{a})$ contains $f$, $V(\mathfrak{a})\cap X_f=\varnothing$ by \hyperref[ex5.23]{\Cref*{ex5.23}(iii)}.
		\item[(iii) $\Rightarrow$ (i)] If not Jacobson, there is a prime ideal $\mathfrak{p}$ such that $\mathfrak{p}\subsetneq\bigcap_{\mathfrak{q}\supseteq\mathfrak{p}}\mathfrak{q}$. Let $f\in\left.\bigcap_{\mathfrak{q}\supseteq\mathfrak{p}}\mathfrak{q}\middle\backslash\mathfrak{p}\right.$, then $X_f\cap V(\mathfrak{p})=\{\mathfrak{p}\}$.
	\end{enumerate}
\end{exercise}

\begin{exercise}\label{ex5.27}
	\begin{enumerate}
		\item[$\Rightarrow$:] Use \Cref{5.21} to proof maximal elements are valuation rings.
		\item[$\Leftarrow$:] If $(A,\mathfrak{m})\subsetneq(A^\prime,\mathfrak{m}^\prime)$, pick $x\in A^\prime\backslash A$. Then $x^{-1}\in\mathfrak{m}$ because it is not a unit in $A$.
	\end{enumerate}
\end{exercise}

\begin{exercise}\label{ex5.28}
	\begin{enumerate}
		\item[(1) $\Rightarrow$ (2)] Pick $a\in\mathfrak{a}\backslash\mathfrak{b}$, and any $b\in\mathfrak{b}$. Consider $a/b$ and $b/a$.
		\item[(2) $\Rightarrow$ (1)] Consider $(a)\subseteq(b)$ or $(b)\subseteq(a)$, we have $a=xb$ or $b=ya$, so $y=x^{-1}$.
	\end{enumerate}
\end{exercise}

\begin{exercise}\label{ex5.29}
	\textcolor{magenta}{$B$ is a local ring of $A$ means $B=A_\mathfrak{p}$ for some prime ideal $\mathfrak{p}$.} Let $\mathfrak{m}$ be the maximal ideal of $B$ and let $\mathfrak{p}=\mathfrak{m}\cap A$. If $s\in A$ but $s\notin\mathfrak{p}$, then $s\in B$ and $s\notin\mathfrak{m}$, so $s$ is a unit in $B$. Let $A_\mathfrak{p}\rightarrow B$ to be $x/s\mapsto s^{-1}x$, which is injective; $\mathfrak{p}A_\mathfrak{p}\subseteq B$, so $B$ dominates $A_\mathfrak{p}$. By \Cref{ex5.28}, $A_\mathfrak{p}$ is a valuation ring, then use \Cref{ex5.27}.
\end{exercise}

\begin{exercise}\label{ex5.30}
	$xy^{-1}\in A$ or $x^{-1}y\in A$, so $\geqslant$ is a total order. Also $(xz)(yz)^{-1}=xy^{-1}$.
	
	\textcolor{magenta}{Treat $v$ as a representation map.} If $v(x)\geqslant v(y)$, $(x+y)y^{-1}=xy^{-1}+1\in A$.
\end{exercise}

\begin{exercise}\label{ex5.31}
	Use the given conditions to verify the set is closed under addition and multiplication. Then verify it is a valuation ring by taking $x=y=1$ and $y=x^{-1}$ in the first condition.
\end{exercise}

\begin{exercise}\label{ex5.32}
	
\end{exercise}

\begin{exercise}\label{ex5.33}
	
\end{exercise}

\begin{exercise}\label{ex5.34}
	
\end{exercise}

\begin{exercise}\label{ex5.35}
	
\end{exercise}

	\section{Chain Conditions}

\begin{exercise}\label{ex6.1}
	\begin{enumerate}[label=(\roman*),ref=\theexercise(\roman*)]\crefalias{enumi}{exercise}
		\item Let $x\in\ker(u)$. Since $u^n$ surjective, there is a $y\in M$ such that $u^n(y)=x$. But $u^{n+1}(y)=u(x)=0$, means that $y\in\ker\left(u^{n+1}\right)=\ker\left(u^n\right)$.
		\item Let $x\in M$. Then $u^n(x)\in\operatorname{im}\left(u^n\right)=\operatorname{im}\left(u^{n+1}\right)$, so there is a $y\in M$ such that $u^{n+1}(y)=u^n(x)$. But $u^n$ injective, means that $u(y)=x$.
	\end{enumerate}
\end{exercise}

\begin{exercise}\label{ex6.2}
	If $M$ is not Noetherian, there is a submodule $N$ which is not finitely generated. Construct a chain of finitely generated submodules by the generators of $N$.
\end{exercise}

\begin{exercise}\label{ex6.3}
	Consider $0\rightarrow N_2/\left(N_1\cap N_2\right)\rightarrow M/\left(N_1\cap N_2\right)\rightarrow M/N_2\rightarrow0$. Use \Cref{2.1}.
\end{exercise}

\begin{exercise}\label{ex6.4}
	Let $m_1,\cdots,m_n$ be generators of $M$, then $\mathfrak{a}=\bigcap\operatorname{Ann}\left(m_i\right)$. $A/\operatorname{Ann}\left(m_i\right)\cong Am_i\subseteq M$, then use \Cref{ex6.3}.
\end{exercise}

\begin{exercise}\label{ex6.5}
	Use the open set definition.
\end{exercise}

\begin{exercise}\label{ex6.6}
	\begin{enumerate}
		\item[(i) $\Rightarrow$ (iii)] Use \Cref{ex6.5}.
		\item[(ii) $\Rightarrow$ (i)] Consider a chain of ascending open sets in $X$.
	\end{enumerate}
\end{exercise}

\begin{exercise}\label{ex6.7}
	Let $\Sigma$ be the collection of closed subsets of $X$ which are not finite union of irreducible closed subspaces. Then $\Sigma$ has minimal element $X_0$. Proof that $X_0$ cannot be either irreducible or reducible.
\end{exercise}

\begin{exercise}\label{ex6.8}
	The converse isn't true. The counter example is $k[x_1,\cdots]/\left(x_1^2,\cdots\right)$.
\end{exercise}

\begin{exercise}\label{ex6.9}
	The irreducible components of $\operatorname{Spec}(A)$ is $V(\mathfrak{p})$ where $\mathfrak{p}$ minimal (\Cref{ex1.20.4}).
\end{exercise}

\begin{exercise}\label{ex6.10}
	Use \Cref{ex3.19.5}.
\end{exercise}

\begin{exercise}\label{ex6.11}
	By \Cref{ex6.7}, $V(\mathfrak{b})$ is a finite union of $V(\mathfrak{q}_i)$, where $\mathfrak{q}_i$ minimal and containing $\mathfrak{b}$. Then use \hyperref[ex5.10]{\Cref*{ex5.10}(c)}.
\end{exercise}

\begin{exercise}\label{ex6.12}
	The converse isn't true. A counter example is an infinite product $\prod\mathbb{Z}/2\mathbb{Z}$.
\end{exercise}
	\section{Noetherian Rings}

\begin{exercise}[Cohen's theorem]\label{ex7.1}
	Let $\mathfrak{a}$ be a maximal element in $\Sigma$. If $xy\in\mathfrak{a}$ but $x,y\notin\mathfrak{a}$, then $\mathfrak{a}+(x)$ is finitely generated. Let $\left\{x,a_1,\cdots,a_n\right\}$ be its generators and let $\mathfrak{a}_0=\left(a_1,\cdots,a_n\right)$. Proof that $\mathfrak{a}_0\nsubseteq\mathfrak{a}$ is impossible. Therefore, if $u\in\mathfrak{a}\subseteq\mathfrak{a}+(x)=\mathfrak{a}_0+(x)$, $u=a_0+bx$ with $b\in\mathfrak{a}$, so $\mathfrak{a}\subseteq\mathfrak{a}_0+x\cdot(\mathfrak{a}:x)$; we also have $x\cdot(\mathfrak{a}:x)\subseteq\mathfrak{a}$. Note that $y\in(\mathfrak{a}:x)$.
\end{exercise}

\begin{exercise}\label{ex7.2}
	\begin{enumerate}
		\item[$\Rightarrow$:] \Cref{ex1.5.2}.
		\item[$\Leftarrow$:] $\left(a_0,a_1,\cdots\right)\subseteq\mathfrak{R}$ and use \Cref{7.15}.
	\end{enumerate}
\end{exercise}

\begin{exercise}\label{ex7.3}
	\begin{enumerate}
		\item[(i) $\Rightarrow$ (ii)] Use \Cref{4.8} and \Cref{4.4}.
		\item[(ii) $\Rightarrow$ (iii)] Let $S=\left\{1,x,x^2,\cdots\right\}$. Proof that $\bigcup\left(\mathfrak{a}:x^n\right)=S(\mathfrak{a})$.
		\item[(iii) $\Rightarrow$ (i)] Similar to the proof of \Cref{7.12}.
	\end{enumerate}
\end{exercise}

\begin{exercise}\label{ex7.4}
	\begin{enumerate}[label=(\roman*),ref=\theexercise(\roman*)]\crefalias{enumi}{exercise}
		\item $\mathfrak{p}=\left\{(z-c):|c|=1\right\}$ is prime, and the ring is $(\mathbb{C}[x])_\mathfrak{p}$. Use \Cref{7.3}.
		\item Use \Cref{ex1.5.1}. If $f=z^ng$ with $g$ invertible, $(f)=\left(z^n\right)$, so the only ideals are $(0)$ and $\left(z^n\right)$.
		\item $\cdots\subsetneq(\sin\pi nz)\subsetneq(\sin\pi(n+1)z)\subsetneq\cdots$ is an ascending chain.
		\item The ring is $\mathbb{C}\left[z^{k+1}\right]$.
		\item The ring is $\mathbb{C}\left[z,z^2,\cdots,wz,wz^2,\cdots,w^2z,w^2z^2,\cdots\right]$.
	\end{enumerate}
\end{exercise}

\begin{exercise}\label{ex7.5}
	Use \Cref{ex5.12} and \Cref{7.8}.
\end{exercise}

\begin{exercise}\label{ex7.6}
	\textcolor{magenta}{A ring is finitely generated if it is finitely generated as a $\mathbb{Z}$-algebra.} $\mathbb{Z}$ is a Jacobson ring. By \Cref{ex5.25}, $K$ is finite over $\mathbb{Z}$, so it is of the form $\mathbb{Z}^n\oplus\bigoplus_i\mathbb{Z}/p_i\mathbb{Z}$ due to the addition operation forms a finitely generated abelian group. If $\operatorname{char} K=0$, $\mathbb{Z}\subseteq\mathbb{Q}\subseteq\mathbb{Z}^n$, then use \Cref{7.8}.
\end{exercise}

\begin{exercise}\label{ex7.7}
	Use ascending chain condition.
\end{exercise}

\begin{exercise}\label{ex7.8}
	$A\cong A[x]/(x)$.
\end{exercise}

\begin{exercise}\label{ex7.9}
	Let $\mathfrak{a}_0\subseteq\mathfrak{a}_1\subseteq\cdots$ be a chain in $A$. Since $\mathfrak{a}_0$ is contained in only finite many maximal ideals, there is a large enough $n$ such that $S_{\mathfrak{m}}^{-1}\mathfrak{a}_n=S_{\mathfrak{m}}^{-1}\mathfrak{a}_{n+1}=\cdots$ for all $\mathfrak{m}\supseteq\mathfrak{a}_0$. We also have $S_{\mathfrak{m}}^{-1}\mathfrak{a}_i=A_\mathfrak{m}$ for all $i$, when $\mathfrak{m}\nsupseteq\mathfrak{a}_0$. Then apply \Cref{3.8} on $S_\mathfrak{m}^{-1}\left(\mathfrak{a}_{n+i}\middle/\mathfrak{a}_n\right)$.
\end{exercise}

\begin{exercise}\label{ex7.10}
	$M[x]=M\otimes A[x]$, so it is a homomorphic image of $M\times A[x]$. Use \Cref{7.1}.
\end{exercise}

\begin{exercise}\label{ex7.11}
	Counter example is an infinite product $\prod\mathbb{Z}/2\mathbb{Z}$.
\end{exercise}

\begin{exercise}\label{ex7.12}
	By \hyperref[ex3.16]{\Cref*{ex3.16}(i)}, $\mathfrak{a}\mapsto\mathfrak{a}^e$ is injective. Then use ascending chain condition.
\end{exercise}

\begin{exercise}\label{ex7.13}
	Transfer the Noetherian property $B\rightarrow B_\mathfrak{p}\rightarrow B_\mathfrak{p}/\mathfrak{p}B_\mathfrak{p}$ and use \Cref{ex6.8}.
\end{exercise}

\begin{exercise}[\textcolor{red}{Hilbert's Nullstellensatz}]\label{ex7.14}
	If $f\notin\sqrt{\mathfrak{a}}$, there is a $\mathfrak{p}\supseteq\mathfrak{a}$ such that $f\notin\mathfrak{p}$. $(A/\mathfrak{p})_{\overline{f}}$ is a finitely generated $k$-algebra, so by \hyperref[7.9]{Zariski's lemma}, $\left.\left((A/\mathfrak{p})_{\overline{f}}\right)\middle/\mathfrak{m}\right.=k$. Let $x_i$ be the images of $t_i$ in $\left.\left((A/\mathfrak{p})_{\overline{f}}\right)\middle/\mathfrak{m}\right.$, then $x=\left(x_1,\cdots,x_n\right)$ satisfies $x\in V$ but $f(x)\neq0$.
	
	\textcolor{red}{Rabinowitsch trick:} \textcolor{green!70!black}{Let $f_1,\cdots,f_m$ be generators of $\mathfrak{a}$. If $f\in I(V)$, then $f_1,\cdots,f_m,1-t_0f$ has no zeros in $k\left[t_0,t_1,\cdots,t_n\right]$. Apply \Cref{ex5.17}, $1=g_0\left(1-t_0f\right)+\sum_{i=1}^mg_if_i$. Now put $t_0=1/f$, we have $g_i\in k\left[t_1,\cdots,t_n,1/f\right]$. Therefore, $g_i=h_i/f^{N_i}$ with $h_i\in k\left[t_1,\cdots,t_n\right]$. Hence we have $f^N=\sum_{i=1}^mh_if^{N-N_i}f_i$ for all $N\geqslant\max_i\left\{N_i\right\}$.}
\end{exercise}

\begin{exercise}\label{ex7.15}
	\begin{enumerate}
		\item[(iv) $\Rightarrow$ (i)] Let $x_1,\cdots,x_n$ be elements of $M$ whose images in $M/\mathfrak{m}M$ form a basis of $k$-vector space. Let $N=\left(x_1,\cdots,x_n\right)$, so $N+\mathfrak{m}M=M$. Since $\mathfrak{J}=\mathfrak{m}$, apply \hyperref[2.6]{Nakayama's lemma} on $M/N$, we have $M=N$. Let $F$ be a free $A$-module with basis $e_1,\cdots,e_n$, and $\phi:F\rightarrow M$ be $\phi\left(e_i\right)=x_i$. Then 
		\begin{equation*}
			\operatorname{Tor}_1^A(k,M)\rightarrow k\otimes_A\ker\phi\rightarrow k\otimes_AF\xrightarrow{1\otimes\phi}k\otimes_AM\rightarrow0
		\end{equation*}
		exact. Since $k\otimes F$ and $k\otimes M=M/\mathfrak{m}M$ are vector spaces of a same dimension, $1\otimes\phi$ is an isomorphism, so $k\otimes\ker\phi=0$. By \Cref{ex2.3}, $\ker\phi=0$.
	\end{enumerate}
\end{exercise}

\begin{exercise}\label{ex7.16}
	Use \Cref{3.10} and \Cref{ex7.15}.
\end{exercise}

\begin{exercise}\label{ex7.17}
	
\end{exercise}

\begin{exercise}\label{ex7.18}
	
\end{exercise}

\begin{exercise}\label{ex7.19}
	
\end{exercise}

\begin{exercise}\label{ex7.20}
	\begin{enumerate}[label=(\roman*),ref=\theexercise(\roman*)]\crefalias{enumi}{exercise}
		\item $\Rightarrow$: We have a chain $\mathcal{F}_0\subseteq\mathcal{C}_1\subseteq\mathcal{F}_1\subseteq\mathcal{C}_2\subseteq\cdots\subseteq\mathcal{F}$, where $\mathcal{F}_0$ is the topology, $\mathcal{C}_i=\mathcal{F}_{i-1}\cup\left\{S^c:S\in\mathcal{F}_{i-1}\right\}$, and $\mathcal{F}_i$ is a collection of all finite intersections of sets in $\mathcal{C}_i$.
		\item $\Rightarrow$: Use the open set definition.
	\end{enumerate}
\end{exercise}

\begin{exercise}\label{ex7.21}
	\begin{enumerate}
		\item[$\Rightarrow$:] Use \Cref{ex7.20}.
		\item[$\Leftarrow$:] If $E\notin\mathcal{F}$, let $\mathcal{C}=\left\{X^\prime\subseteq X:X^\prime\text{ closed, }E\cap X^\prime\notin\mathcal{F}\right\}$. By Noetherian property, $\mathcal{C}$ has a minimal element $X_0$. Use the closed set definition to show $X_0$ is irreducible. If $\overline{E\cap X_0}\neq X_0$, $E\cap\overline{E\cap X_0}\in\mathcal{F}$ by minimality, but $E\cap\overline{E\cap X_0}=E\cap X_0$. If $U_0\subseteq E\cap X_0$, let $U_0=U\cap X_0$ where $U$ is open in $X$, then $E\cap X_0=U_0\cup\left(E\cap U^c\cap X_0\right)$.
	\end{enumerate}
\end{exercise}

\begin{exercise}\label{ex7.22}
	\begin{enumerate}
	\item[$\Leftarrow$:] If $E$ is not open, let $\mathcal{C}=\left\{X^\prime\subseteq X:X^\prime\text{ closed, }E\cap X^\prime\text{ is not open}\right\}$ and use the notations similar to \Cref{ex7.21}. If $U_0\subseteq E\cap X_0$, $U^c\cap X_0$ is closed, so by minimality $E\cap U^c\cap X_0$ is open.
	\end{enumerate}
\end{exercise}

\begin{exercise}\label{ex7.23}
	Hint.
\end{exercise}

\begin{exercise}\label{ex7.24}
	Hint.
\end{exercise}

\begin{exercise}\label{ex7.25}
	Use \Cref{ex5.11} and \Cref{ex7.24}.
\end{exercise}

\begin{exercise}\label{ex7.26}
	
\end{exercise}

\begin{exercise}\label{ex7.27}
	
\end{exercise}
	\section{Artinian Rings}

\begin{exercise}\label{ex8.1}
	Use \Cref{7.16} on local ring $A_{\mathfrak{p}_i}$. Then apply the contraction. If $\mathfrak{q}_i$ is isolated, then $\mathfrak{p}_i$ is minimal.
\end{exercise}

\begin{exercise}\label{ex8.2}
	\begin{enumerate}
		\item[(i) $\Rightarrow$ (ii)] Use \Cref{8.1}.
		\item[(iii) $\Rightarrow$ (i)] Use \Cref{8.5}.
	\end{enumerate}
\end{exercise}

\begin{exercise}\label{ex8.3}
	\begin{enumerate}
		\item[(i) $\Rightarrow$ (ii)] By \Cref{8.7}, we can assume that $A$ is an Artinian local ring. By \Cref {8.4}, $\mathfrak{m}^n=0$ for some $n$, so we have a chain $A=\mathfrak{m}^0\supseteq\mathfrak{m}^1\supseteq\cdots\supseteq\mathfrak{m}^n=0$.\par
		By \hyperref[7.9]{Zariski's lemma}, $A/\mathfrak{m}$ is a finitely generated $k$-module. By \Cref{8.5}, $A$ is Noetherian, so $\mathfrak{m}^i$ is a finitely generated $A$-module,  and therefore $\mathfrak{m}^i/\mathfrak{m}^{i+1}$ is a finitely generated $A/\mathfrak{m}$-module. Use
		\begin{align*}
			&\dim_k\left(\mathfrak{m}^i\middle/\mathfrak{m}^{i+1}\right)=\dim_{A/\mathfrak{m}}\left(\mathfrak{m}^i\middle/\mathfrak{m}^{i+1}\right)\dim_k\left(A\middle/\mathfrak{m}\right);&&\\
			&\dim_k\left(\mathfrak{m}^i\right)=\dim_k\left(\mathfrak{m}^{i+1}\right)+\dim_k\left(\mathfrak{m}^i\middle/\mathfrak{m}^{i+1}\right),&&\text{where }0\rightarrow\mathfrak{m}^{i+1}\rightarrow\mathfrak{m}^i\rightarrow\mathfrak{m}^i/\mathfrak{m}^{i+1}\rightarrow0
		\end{align*}
		to finish the proof.
		\item[(ii) $\Rightarrow$ (i)] $A$ is a $k$-vector space of finite dimension. Use \Cref{6.10}.
	\end{enumerate}
\end{exercise}

\begin{exercise}\label{ex8.4}
	\begin{enumerate}
		\item[(i) $\Rightarrow$ (iii)] The finitely-generated property is preserved for localization and quotient of modules.
		\item[(ii) $\Rightarrow$ (iii)] $B_\mathfrak{p}/\mathfrak{p}B_\mathfrak{p}$ is a finitely generated $k(\mathfrak{p})$-algebra. By \hyperref[7.5]{Hilbert's basis theorem}, any finitely generated $k(\mathfrak{p})$-algebra is Noetherian. Then use \Cref{ex8.2} and \Cref{ex8.3}.
		\item[(iii) $\Rightarrow$ (ii),(iv)] $B_\mathfrak{p}/\mathfrak{p}B_\mathfrak{p}$ is Noetherian and integral over the field $k(\mathfrak{p})$. If $\mathfrak{q}\in B_\mathfrak{p}/\mathfrak{p}B_\mathfrak{p}$ is prime, $\left(B_\mathfrak{p}/\mathfrak{p}B_\mathfrak{p}\right)/\mathfrak{q}$ is an integral domain, so it is a field by \Cref{5.7}. Hence, every prime ideal in $B_\mathfrak{p}/\mathfrak{p}B_\mathfrak{p}$ is maximal, which means $\operatorname{dim}\left(B_\mathfrak{p}/\mathfrak{p}B_\mathfrak{p}\right)=0$. Thus, $B_\mathfrak{p}/\mathfrak{p}B_\mathfrak{p}$ is Artinian by \Cref{8.5}. Then use \Cref{8.3}. 
	\end{enumerate}
	A counter example is $f:\mathbb{Q}\hookrightarrow\overline{\mathbb{Q}}$ where $\overline{\mathbb{Q}}$ is an algebraic closure of $\mathbb{Q}$.
\end{exercise}

\begin{exercise}\label{ex8.5}
	Use \hyperref[noethernormalization]{Noether normalization lemma} and let $\mathfrak{m}$ be a maximal ideal in $k\left[y_1,\cdots,y_r\right]$. By \Cref{ex8.4}, $A_\mathfrak{m}/\mathfrak{m}A_\mathfrak{m}$ is finite over $k\left[y_1,\cdots,y_r\right]/\mathfrak{m}=k$, so by \Cref{ex8.3} and \Cref{8.3}, $A_\mathfrak{m}/\mathfrak{m}A_\mathfrak{m}$ only has finite many maximal ideals, correspond to finite many points in $X$.
	
	Let $N$ be the number of generators of $A$ over $k\left[y_1,\cdots,y_r\right]$. The $k$-dimension of $k$-vector space $A_\mathfrak{m}/\mathfrak{m}A_\mathfrak{m}$ is uniformly bounded by $N$. If $\mathfrak{n}_1,\cdots\mathfrak{n}_M$ are maximal ideals in $A_\mathfrak{m}/\mathfrak{m}A_\mathfrak{m}$, by \Cref{1.10}, $A_\mathfrak{m}/\mathfrak{m}A_\mathfrak{m}\rightarrow\prod_{i=1}^M\left.\left(A_\mathfrak{m}/\mathfrak{m}A_\mathfrak{m}\right)\middle/\mathfrak{n}_i\right.$ is surjective, so $\dim_k\left(A_\mathfrak{m}/\mathfrak{m}A_\mathfrak{m}\right)\geqslant\sum_{i=1}^M\dim_k\left(\left.\left(A_\mathfrak{m}/\mathfrak{m}A_\mathfrak{m}\right)\middle/\mathfrak{n}_i\right.\right)\geqslant M$.
\end{exercise}

\begin{exercise}\label{ex8.6}
 	Consider $(A/\mathfrak{q})_{\overline{\mathfrak{p}}}$, which is an Artinian local ring. Use the following claims then \Cref{6.7}.
 	\begin{itemize}
 		\item \textit{If $A$ is an Artinian local ring, $\mathfrak{m}$ its only maximal ideal. Then every ideal in $A$ is $\mathfrak{m}$-primary.}\\ Radical is the intersection of the only prime ideal.
 		\item \textit{There is a 1-1 correspondence between $\mathfrak{p}$-primary ideals in $A$ and $S^{-1}\mathfrak{p}$-primary ideals in $S^{-1}A$.}\\
 		Use \Cref{3.11.1} and \Cref{4.9}.
 	\end{itemize}
\end{exercise}
	\section{Discrete Valuation Rings and Dedekind Domain}

\begin{exercise}\label{ex9.1}
	\Cref{3.11.4} asserts that $\dim S^{-1}A$ can only be 0 and 1. If $\dim S^{-1}A=0$, it is $\operatorname{Frac}A$; if $\dim S^{-1}A=1$, use \Cref{5.12}.
	
	By \Cref{3.15}, $S^{-1}\left(M^{-1}\right)=S^{-1}(A:M)=\left(S^{-1}A:S^{-1}M\right)=\left(S^{-1}M\right)^{-1}$. We also have $S^{-1}(MN)=\left(S^{-1}M\right)\left(S^{-1}N\right)$, so the map $M\mapsto S^{-1}M$ is a homomorphism. Use \Cref{3.11.1}. Note that $(u)\mapsto S^{-1}(u)=\left(S^{-1}A\right)u$ maps $P$ into $P^\prime$, so $H\rightarrow H^\prime$ is surjective.
\end{exercise}

\begin{exercise}\label{ex9.2}
	Localize and use \Cref{3.8} to the final move, we can assume that $A$ is a DVR.
	
	\hyperref[9.2]{\Cref*{9.2}(iii)(v)} assert that $c(f)=\left(x^n\right)$ where $x$ is the generator of $\mathfrak{m}$, so $f=x^nf_0$ where $f_0$ is primitive. Apply \Cref{ex1.2.4} to $f_0$, $g_0$ and note that for any $a\in A$, $ac(f)=c(af)$.
\end{exercise}

\begin{exercise}\label{ex9.3}
	\begin{enumerate}
		\item[$\Rightarrow$:] $\left(x_1,x_2\right)=\left(x_1\right)$ if $x_1x_2^{-1}\in A$, so every ideal is principal. Then use \Cref{9.2.3}.
		\item[$\Leftarrow$:] If $v(y)\geqslant v(x)$, $y=\left(yx^{-1}\right)x$, so every ideal in DVR is of the form $\mathfrak{m}_k=\left\{z\in A:v(z)\geqslant k\right\}$.
	\end{enumerate}
\end{exercise}

\begin{exercise}\label{ex9.4}
	For any non-zero $x\in A$, $x\in\mathfrak{m}^n$ but $x\notin\mathfrak{m}^{n+1}$ for some $n$. Define a discrete valuation $v(x)=n$.
\end{exercise}

\begin{exercise}\label{ex9.5}
	By \Cref{ex3.13} and \Cref{ex7.16}, and notice that torsion free module over a PID is flat.
\end{exercise}

\begin{exercise}\label{ex9.6}
	From \Cref{9.3.3}, \Cref{9.2.5}, and the structure theorem for finitely generated modules over a PID, $M_\mathfrak{p}=\bigoplus_{i=1}^NA_\mathfrak{p}/\mathfrak{p}^{n_i}A_\mathfrak{p}=\bigoplus_{i=1}^N\left(A/\mathfrak{p}^{n_i}\right)_\mathfrak{p}$. Define a homomorphism $\phi:M\rightarrow\bigoplus_\mathfrak{p}M_\mathfrak{p}$ by $m\mapsto(m/1,\cdots)$, so each $\phi_\mathfrak{p}$ is an isomorphism. By \Cref{3.9}, $\phi$ is an isomorphism.
	
	By \Cref{ex3.19.5}, $\operatorname{Supp}(M)=V(\operatorname{Ann}(M))$. From \Cref{9.1} and \Cref{9.3.2}, $\operatorname{Ann}(M)=\mathfrak{p}_1^{n_1}\cdots\mathfrak{p}_k^{n_k}$. Use $\dim A=1$ to proof that $V(\operatorname{Ann}(M))$ is finite.
\end{exercise}

\begin{exercise}\label{ex9.7}
	From \Cref{9.1} and \Cref{9.3.2}, $\mathfrak{a}=\mathfrak{p}_1^{n_1}\cdots\mathfrak{p}_k^{n_k}$. By \Cref{1.10}, $A\rightarrow\prod_{i=1}^{k}A/\mathfrak{p}_i^{n_i}$ is a surjective with kernel $\mathfrak{a}$. Note that each component $A/\mathfrak{p}_i^{n_i}$ is local so is PID.
	
	If $a\in\mathfrak{b}$, $\mathfrak{b}/(a)=(\overline{b})$ and therefore $\mathfrak{b}=(a,b)$.
\end{exercise}

\begin{exercise}\label{ex9.8}
	Localize and note that all ideals in DVR are $\left\{z\in A:v(z)\geqslant k\right\}$. Use \Cref{3.8} to pull back.
\end{exercise}

\begin{exercise}\label{ex9.9}
	This is equivalent to say that
	\begin{equation*}
		A\xrightarrow{\phi}\bigoplus_{i=1}^nA/\mathfrak{a}_i\xrightarrow{\psi}\bigoplus_{i<j}A/\left(\mathfrak{a}_i+\mathfrak{a}_j\right)
	\end{equation*}
	is exact, where $\phi(x)=\left(x+\mathfrak{a}_i\right)_{i=1}^n$, and $\psi\left(x_i+\mathfrak{a}_i\right)_{i=1}^n=\left(x_i-x_j+\mathfrak{a}_i+\mathfrak{a}_j\right)_{i<j}$. If $A$ is DVR, we have $\mathfrak{a}_i+\mathfrak{a}_j=\mathfrak{a}_{\min(i,j)}$. If $\left(x_i+\mathfrak{a}_i\right)_{i=1}^n\in\ker\psi$, $x_i+\mathfrak{a}_i=x_j+\mathfrak{a}_i$ for all $i<j$, so take $j=n$ for all $i$ and we have $\phi\left(x_n\right)\in\ker\psi$. Then use \Cref{9.3.3} and \Cref{3.8} for Dedekind domain.
\end{exercise}
	\newpage
	% Summaries section
	\phantomsection
	\section*{SUMMARY}
	\addcontentsline{toc}{chapter}{SUMMARY}
	\markboth{SUMMARY}{SUMMARY}
	\RedeclareSectionCommand[beforeskip=10pt]{subsubsection}
	\subsubsection*{Chapter 1 - Sheaves}

\textcolor{cyan}{Atiyah-Macdonald Exercise 2.15} is used thoroughly. \textcolor{cyan}{(i)} means $t\in\mathcal{F}_x$ can be represented by $s\in\mathcal{F}(U)$, \textcolor{cyan}{(ii)} is used for finding agreement on a smaller neighborhood.

The restriction $\mathcal{F}|_Z$ can be simplified if $Z$ is an open set.

\subsubsection*{Chapter 2 - Schemes}

Most exercises help to enhance the comprehension of the definitions.

\subsubsection*{Chapter 3 - First Properties of Schemes}

The definitions generalize the results in commutative algebra. Gluing morphisms and schemes is a new technique and is used through the exercises.

There is an important theorem which is proved in \Cref{ex3.8}: The intersection of two affine subschemes can be covered by $\left\{W_i\right\}$, where each $W_i$ is basic open in both subschemes.

\Cref{ex3.11.2} is useful when deal with closed subschemes.

\subsubsection*{Chapter 4 - Separated and Proper Morphisms}

\Cref{4.4} and \Cref{4.8} are fundamental for understanding primary ideal structure.

\Cref{4.9} is useful for computations.

\subsubsection*{Chapter 5 - Integral Dependence and Valuations}

The most important chapter in the book.

\textbf{Important results:} \hyperref[5.10]{Lying Over Theorem}, \hyperref[noethernormalization]{Noether Normalization Lemma}.

\Cref{ex5.10} connects the going-up and going-down properties of a ring homomorphism with topological properties of the induced map on spectra, but the ideas are not used further in the book.

\Cref{ex5.25} uses Jacobson rings to generalize Zariski's Lemma.

\subsubsection*{Chapter 6 \& 7 - Chain Conditions \& Noetherian Rings}

Most exercises further refine old concepts rather than introducing new ideas.

\subsubsection*{Chapter 8 - Artinian Rings}

An important result is \Cref{8.5}.

\Cref{ex8.3} provides another generalization of Zariski's Lemma.

\subsubsection*{Chapter 9 - Discrete Valuation Rings and Dedekind Domains}

This chapter is closely connected to algebraic number theory. \Cref{9.2} and \Cref{9.3} appear repeatedly throughout the exercises.

\end{document}
